% ClaimFlow Sprint 2 report (CPE 334 Deliverable 2). Build: bun run pdf 02
% in docs/reports, or save in VS Code with LaTeX Workshop.
\documentclass[11pt]{article}
\newcommand{\sprintname}{Sprint 2}
\newcommand{\deliverable}{Project Deliverable 2}
% ---------------------------------------------------------------------------
% Shared preamble of every ClaimFlow sprint report. Build with XeLaTeX from
% the sprint folder (bun run pdf sprint-N). Before \input{../preamble},
% main.tex defines \sprintname (Sprint 1) and \deliverable (Project
% Deliverable 1).
%
% Fonts: Times New Roman and Courier New ship with Windows and macOS and both
% carry the arrows and comparison signs the sources use. A glyph neither has
% shows as "Missing character" in main.log; check-report counts those.
% ---------------------------------------------------------------------------
\usepackage{fontspec}
\setmainfont{Times New Roman}
\setmonofont{Courier New}

% Student names and IDs, from docs/reports/.env through bun run identity.
\InputIfFileExists{build/identity.tex}{}{%
  \errmessage{build/identity.tex is missing: run bun run identity <sprint>}}

\usepackage[a4paper, margin=2cm]{geometry}
\usepackage{graphicx}
\usepackage{pdflscape}
\usepackage{booktabs}
\usepackage{longtable}
\usepackage{array}
\usepackage{xcolor}
\usepackage{fvextra}
\usepackage{float}
\usepackage{caption}
\usepackage{etoolbox}
\usepackage{refcount}
\usepackage{seqsplit}
\usepackage{fancyhdr}

% Colours before hyperref, which reads linkcolor as it loads.
\definecolor{brand}{HTML}{1B5E3F}
\definecolor{todored}{HTML}{C62828}
\usepackage[colorlinks=true, urlcolor=blue!55!black, linkcolor=brand,
            pdftitle={ClaimFlow \sprintname\space report},
            pdfauthor={\studentOneName; \studentTwoName}]{hyperref}

\setlength{\parindent}{0pt}
\setlength{\parskip}{5pt}
\setlength{\emergencystretch}{3em}
\hfuzz=3pt % ignore overfull boxes under 3 pt, which no reader can see
\captionsetup{font=small}
\setcounter{tocdepth}{2}
\setcounter{secnumdepth}{0}

\pagestyle{fancy}
\fancyhf{}
\fancyfoot[L]{\footnotesize ClaimFlow, \sprintname\space report}
\fancyfoot[R]{\footnotesize page \thepage}
\renewcommand{\headrulewidth}{0pt}
\fancypagestyle{plain}{\fancyhf{}\fancyfoot[R]{\footnotesize page \thepage}%
  \renewcommand{\headrulewidth}{0pt}}

% ---- Evidence that is not there yet ---------------------------------------

% Red inline text. check-report counts every \todo and \pending.
\newcommand{\todo}[1]{\textcolor{todored}{\textbf{[TODO: #1]}}}

% A red box where a capture or section is still waiting for something.
% \pending{what it shows}{what it waits for}
\newcommand{\pending}[2]{%
  \par\medskip\noindent\fcolorbox{todored}{white}{\parbox{0.97\linewidth}{%
    \textcolor{todored}{\textbf{[PENDING: #1]}}\\
    \small\itshape Waits for: #2}}\par\medskip}

% ---- Parts -----------------------------------------------------------------

% \reportpart{A}{Title}: starts a page, adds the contents entry, and sets
% part:A:start. \reportpartend{A} sets part:A:end on the last page, so the
% navigation table can print real page ranges with \pageref.
\newcommand{\reportpart}[2]{%
  \clearpage
  \phantomsection
  \addcontentsline{toc}{section}{Part #1: #2}%
  \label{part:#1:start}%
  {\noindent\bfseries\Large\color{brand} Part #1: #2\par}%
  \vspace{2pt}\noindent\textcolor{brand}{\rule{\linewidth}{1.2pt}}\par
  \vspace{6pt}}
\newcommand{\reportpartend}[1]{\par\label{part:#1:end}}
% \pages{A} prints "pp. 5--12" or "p. 5"
\newcommand{\pages}[1]{%
  \ifnumequal{\getpagerefnumber{part:#1:start}}{\getpagerefnumber{part:#1:end}}%
    {p.\ \pageref{part:#1:start}}%
    {pp.\ \pageref{part:#1:start}--\pageref{part:#1:end}}}

% Sub-headings inside a part.
\newcommand{\partsection}[2]{%
  \phantomsection\addcontentsline{toc}{subsection}{#1\ \ #2}%
  \par\bigskip{\noindent\bfseries\large\color{brand} #1\ \ #2\par}\smallskip}

% ---- Evidence figures --------------------------------------------------------

\newcommand{\evidencemaxheight}{22cm}

% The red box drawn when a figure file does not exist, so the build never
% fails on a missing capture.
\newcommand{\figmissing}[2]{%
  \fbox{\parbox[c][3cm][c]{0.94\linewidth}{\centering
    \textcolor{todored}{\textbf{[PENDING FIGURE: \detokenize{#1}]}}\\[4pt]
    \small\itshape #2}}}

% \evidence{figures/x.png}{Caption}: one full-width numbered figure.
\newcommand{\evidence}[2]{%
  \IfFileExists{#1}%
    {\begin{figure}[H]\centering
       \fbox{\includegraphics[width=0.96\linewidth,height=\evidencemaxheight,keepaspectratio]{#1}}%
       \caption{#2}\end{figure}}%
    {\begin{figure}[H]\centering\figmissing{#1}{#2}\caption{#2}\end{figure}}}

% \evidencepart{figures/x-p2.png}{Caption}: one unnumbered piece of a sliced capture.
\newcommand{\evidencepart}[2]{%
  \IfFileExists{#1}%
    {\begin{figure}[H]\centering
       \fbox{\includegraphics[width=0.96\linewidth,height=\evidencemaxheight,keepaspectratio]{#1}}%
       \caption*{\small\itshape #2}\end{figure}}%
    {\begin{figure}[H]\centering\figmissing{#1}{#2}\caption*{\small\itshape #2}\end{figure}}}

% \evidencehalf{file}{caption}: wide capture, at most 11.5 cm tall, two to a page.
\newcommand{\evidencehalf}[2]{%
  \IfFileExists{#1}%
    {\begin{figure}[H]\centering
       \fbox{\includegraphics[width=0.96\linewidth,height=11.5cm,keepaspectratio]{#1}}%
       \caption{#2}\end{figure}}%
    {\begin{figure}[H]\centering\figmissing{#1}{#2}\caption{#2}\end{figure}}}

% \diagramfull{file}{caption}: one Mermaid diagram of a captured document,
% photographed on its own at 3x by scripts/capture-diagrams.ts, at full text width.
\newcommand{\diagramfull}[2]{%
  \IfFileExists{#1}%
    {\begin{figure}[H]\centering
       \fbox{\includegraphics[width=0.96\linewidth,height=\evidencemaxheight,keepaspectratio]{#1}}%
       \caption{#2}\end{figure}}%
    {\begin{figure}[H]\centering\figmissing{#1}{#2}\caption{#2}\end{figure}}}

% \diagramwide{file}{caption}: a diagram much wider than tall, on its own landscape page.
\newcommand{\diagramwide}[2]{%
  \IfFileExists{#1}%
    {\begin{landscape}\begin{figure}[H]\centering
       \fbox{\includegraphics[width=0.96\linewidth,height=15cm,keepaspectratio]{#1}}%
       \caption{#2}\end{figure}\end{landscape}}%
    {\begin{figure}[H]\centering\figmissing{#1}{#2}\caption{#2}\end{figure}}}

% \evidencesmall{file}{caption}{width fraction}: a capture that is small to begin with.
\newcommand{\evidencesmall}[3]{%
  \IfFileExists{#1}%
    {\begin{figure}[H]\centering
       \fbox{\includegraphics[width=#3\linewidth,keepaspectratio]{#1}}%
       \caption{#2}\end{figure}}%
    {\begin{figure}[H]\centering\figmissing{#1}{#2}\caption{#2}\end{figure}}}

% \docsource{docs/srs.md}{sprint-1}{abc1234}: the line above each captured document.
\newcommand{\docsource}[3]{%
  \par\noindent{\small\textbf{Source:} \texttt{\detokenize{#1}} as GitHub renders it,
  captured from \texttt{\detokenize{#2}} at commit \texttt{\detokenize{#3}}.}\par\smallskip}

% ---- Code and terminal output ------------------------------------------------

\DefineVerbatimEnvironment{shell}{Verbatim}
  {fontsize=\footnotesize, frame=single, framesep=5pt, rulecolor=\color{black!25},
   breaklines=true, breaksymbolleft={}}
\DefineVerbatimEnvironment{codeblock}{Verbatim}
  {fontsize=\scriptsize, frame=single, framesep=5pt, rulecolor=\color{black!25},
   breaklines=true, breaksymbolleft={}}

% ---- Title block --------------------------------------------------------------

\newcommand{\reporttitle}{%
  \begin{center}
    {\Large\bfseries CPE 334 --- Introduction to Software Engineering\par}
    \vspace{3pt}
    {\large \deliverable: \sprintname\space report\par}
    \vspace{10pt}
    {\Huge\bfseries\color{brand} ClaimFlow\par}
    {\large Theme 5, Department Expense Reimbursement System\par}
  \end{center}}

\emergencystretch=2em

% The back matter is not a rubric part, so it gets its own heading and labels.
\newcommand{\reportback}[1]{%
  \clearpage\phantomsection\addcontentsline{toc}{section}{#1}%
  \label{part:Z:start}%
  {\noindent\bfseries\Large\color{brand} #1\par}%
  \vspace{2pt}\noindent\textcolor{brand}{\rule{\linewidth}{1.2pt}}\par\vspace{6pt}}
\newcommand{\reportbackend}{\par\label{part:Z:end}}

\begin{document}

\thispagestyle{plain}
\reporttitle
\draftbanner

\medskip
\begin{center}
\begin{tabular}{@{}lll@{}}
\toprule
 & \textbf{Student 1} & \textbf{Student 2} \\ \midrule
Name & \studentOneName & \studentTwoName \\
Student ID & \studentOneId & \studentTwoId \\
GitHub & \href{https://github.com/Kiatisakk}{@Kiatisakk} & \href{https://github.com/beambeambeam}{@beambeambeam} \\
\bottomrule
\end{tabular}
\end{center}

\begin{tabular}{@{}p{3.9cm}p{12.4cm}@{}}
\textbf{Theme} & Theme 5, ClaimFlow: Department Expense Reimbursement System \\
\textbf{Repository} & \url{https://github.com/beambeambeam/claimflow} \\
\textbf{Submitted tag} & \texttt{v0.1.1.0} (semantic-release on \texttt{main}; \url{https://github.com/beambeambeam/claimflow/releases/tag/v0.1.1.0}), published 2026-11-25 12:37 UTC \\
\textbf{Submitted commit} & \texttt{3a91ea4} on \texttt{main}: merge commit of release PR \#109 (\texttt{sprint-2} into \texttt{main}), 2026-11-25 12:33 UTC; CI run 45614889848 and Deploy demo run 45615736377 \\
 & There is no second release: the submitted candidate is \texttt{v0.1.1.0}, the accepted Sprint 2 release, and it contains no Sprint 3 feature (Part C). \\
\textbf{Deployed application} & \url{https://claimflow-5bcxewbqkq-as.a.run.app} (Cloud Run, region \texttt{asia-southeast1}) \\
\textbf{Deployed version} & \texttt{GET /api/v1/health} on this URL answered \texttt{\{"status":"ok",\allowbreak "version":\allowbreak "3a91ea458862a02fc88384f09ef74b7a502ce762"\}} on 2026-11-25 at about 20:20 (UTC+7), right after the release deploy (Deploy demo run 45615736377). The live site now runs a later commit: \texttt{GET /api/v1/health} answered \texttt{\{"status":"ok",\allowbreak "version":\allowbreak "1ee2563de7cddf237e5066e5bdf1c577ae66e9bf"\}} on 2026-11-26 at 00:44 (UTC+7), deployed by Deploy demo run 45656061391 for PR \#112. Commit \texttt{1ee2563} changes one file, \path{docs/reports/labs/07/content/00-front.tex} (the text of the Lab 7 report): the application code and the \texttt{v0.1.1.0} tree are unchanged (\texttt{git diff --stat 3a91ea4 1ee2563}: 1 file, 15 insertions). Details in E.1. \\
\textbf{Evidence source} & Sprint 2 evidence (Part B) is captured from \texttt{v0.1.1.0} by design. Front matter, the Part A demonstration and Part E come from the submitted candidate; the candidate is \texttt{v0.1.1.0} itself, so every source in this report is \texttt{3a91ea4}, except the Sprint 3 contracts (Part C: the DoD at \texttt{ef38550} on \texttt{sprint-3}, the three feature contracts at \texttt{21a317c}, the head of the open diagram fix PR \#114) and the live check of E.1. Exception stated once: the demonstrations were run on a local stack built from the \texttt{3a91ea4} tree, not on the deployed service. \\
\end{tabular}

\medskip
\noindent\textbf{How to read this PDF.}
Parts A to E follow the Deliverable 2 rubric. Contracts, the SRS, the feature inventory, the System-Level SDS, the ADR and the traceability matrix are the repository documents themselves, captured as GitHub renders them from \texttt{3a91ea4}, each headed by its file and commit.
Tables, demonstrations, test output and process evidence are written for this report from commands that were run; every number, commit, pull request and date comes from \texttt{git}, \texttt{gh} or a test run.
Open items are red \textbf{[TODO: \ldots]} markers; the DRAFT box above stays while any exist.
The report is built from \path{docs/reports/sprints/02}; \path{docs/reports/sprints/02/README.md} indexes the evidence.

\clearpage
\section*{Navigation}
\addcontentsline{toc}{section}{Navigation table}
\renewcommand{\arraystretch}{1.25}
\begin{longtable}{@{}p{1.2cm}p{2.3cm}p{12.6cm}@{}}
\toprule \textbf{Part} & \textbf{Pages} & \textbf{Content and main IDs} \\ \midrule \endhead
A & \pages{A} & Sprint 1 continuity and change status: accepted Sprint 1 release \texttt{v0.0.3.0} / \texttt{ac5bce5}; continuity demonstration of the Sprint 1 journey (F02-AC-01, F04-AC-04, F03-AC-05, BR-12); change and defect register (\#86, \#87, \#99, \#66, \#67, \#70, \#61); approved baseline revisions (SRS 0.3, F-04 0.6 to 0.9, Product DoD 0.3, SDS 0.12 to 0.18, plan). \\
B & \pages{B} & Sprint 2 contracts and working increment: F-05 (F05-AC-01 to 20, F05-TEST-01 to 26), F-06 (F06-AC-01 to 13, F06-TEST-01 to 17), F-07 (F07-AC-01 to 10, F07-TEST-01 to 13), Sprint 2 DoD; completion table of 43 AC; demonstrations of F-05, F-06, F-07 with a 403; test excerpts \texttt{F06-TEST-07} (BR-14) and \texttt{F07-TEST-04} (access); TDD red and green for BR-14 and BR-16; Issue and PR chains \#43 to \#45 / PR \#108 and \#39 to \#42 / PR \#104, through \texttt{sprint-2} and PR \#109 to \texttt{main}. \\
C & \pages{C} & Sprint 3 contracts and working progress: F-08, F-09, F-10, the Sprint 3 DoD (PR \#111 at \texttt{ef38550}; diagram fix PR \#114 at \texttt{21a317c}), implemented progress (none; PR \#113 for F-08 started), status and remaining-work tables with owners and dates to 13 December. \\
D & \pages{D} & Updated requirements and test traceability: SRS (FR-01 to FR-14, BR-01 to BR-17, NFR-01 to NFR-10), feature inventory F-01 to F-10, System-Level SDS, ADR 0001, traceability matrix of 41 requirements, coverage gaps. \\
E & \pages{E} & Integrated release candidate: build, test, CI and deployment output with counts, regression and post-deployment checks, setup summary, readiness review, dated plan to 13 December. \\
Back & \pages{Z} & Back matter: contribution table, AI-use declaration, known limitations, failed and skipped tests. \\
\bottomrule
\end{longtable}

\tableofcontents


\reportpart{A}{Sprint 1 continuity and change status}
\partsection{A.1}{Accepted Sprint 1 release and continuity of its journey}\label{A:1}

\textbf{Accepted Sprint 1 release.} Tag \texttt{v0.0.3.0}, commit \texttt{ac5bce5} on \texttt{main}: the merge commit of release PR \#69 (\texttt{sprint-1} at \texttt{7e787d4} into \texttt{main}), 2026-11-15 15:31 (UTC+7). It is the version the Deliverable 1 report cites. \texttt{main} has moved since; the releases between that commit and the Sprint 2 release, from \texttt{git tag} and \texttt{git log}:

\begin{center}\small
\begin{tabular}{@{}llllp{7.6cm}@{}}
\toprule \textbf{Tag} & \textbf{Commit} & \textbf{Created (UTC+7)} & \textbf{Release PR} & \textbf{What it added to Sprint 1} \\ \midrule
\texttt{v0.0.3.0} & \texttt{ac5bce5} & 2026-11-15 15:31 & \#69 & Accepted Sprint 1 release (F-01 to F-04) \\
\texttt{v0.0.4.0} & \texttt{7952445} & 2026-11-18 20:03 & \#76 & Lab 5: cloud-ready image, compose stack, GCP setup script \\
\texttt{v0.1.0.0} & \texttt{c32e645} & 2026-11-25 00:48 & \#106 & Lab 5 deploy pipeline, Lab 6 fixes (\#86, \#87), Sprint 1 follow-ups (\#97), Lab 7 change L7/01 (\#99) \\
\texttt{v0.1.1.0} & \texttt{3a91ea4} & 2026-11-25 19:33 & \#109 & Sprint 2: F-05, F-06, F-07 (accepted Sprint 2 release) \\
\bottomrule
\end{tabular}
\end{center}

\noindent The Lab 6 release PR \#93 (\texttt{bb5a0d1}, deployed 2026-11-21 as Cloud Run revision \texttt{claimflow-00010-985}) produced no tag of its own: semantic-release found no relevant change, and its commits first appear under \texttt{v0.1.0.0}.

\textbf{Continuity demonstration.} The original Sprint 1 journey (submit, review, approve, read back) and one Sprint 1 rule (BR-12), run on the Sprint 2 release. The submitted candidate is \texttt{v0.1.1.0} until the Sprint 3 release exists. \todo{recapture on candidate: rerun \texttt{bun run capture-demo 02} on the stack built from the candidate tag and replace figures \texttt{a1-*}} Demonstrations are taken by \path{docs/reports/sprints/02/demo.ts} (Playwright) on a local stack: PostgreSQL 17 started empty, \texttt{migrate up}, \texttt{seed} (7 users and 9 sample claims; a second \texttt{seed} added 0), the API built from the \texttt{3a91ea4} tree with \texttt{APP\_VERSION} set to that SHA, and the web dev server, on 2026-11-25 in Chrome with synthetic data. The seed password was a throwaway value set only in the process environment; it is in no file and in no figure. Claim numbers in the screenshots (\texttt{CLM-2026-000001} and following) come from that fresh database; a rerun gives the next free numbers.

\newcommand{\demofacts}[7]{%
  \begin{center}\small
  \begin{tabular}{@{}>{\bfseries}p{3.1cm}p{13.2cm}@{}}
  \toprule
  IDs & #1 \\
  Role & #2 \\
  Starting conditions & #3 \\
  Steps & #4 \\
  Expected result & #5 \\
  Actual result & #6 \\
  Backend evidence & #7 \\
  \bottomrule
  \end{tabular}
  \end{center}}

\subsubsection*{Continuity demo: submit, approve, read back (Sprint 1 journey)}

\demofacts
 {F-02, F-03, F-04; F02-AC-01, F04-AC-04, F03-AC-05; BR-12, F02-AC-06; tag \texttt{v0.1.1.0}, commit \texttt{3a91ea4}}
 {Requester \texttt{req.anan}, then reviewer \texttt{rev.malee}, then \texttt{req.anan} again}
 {Fresh seeded database, API and web running}
 {\texttt{req.anan} submits a claim with a Travel item of 250.00 and a Meals item of 400.00; \texttt{rev.malee} opens it from the Review Queue and approves; \texttt{req.anan} opens My Claims and the claim. Then \texttt{req.anan} enters a Meals item of 1,600.00 THB and presses Submit.}
 {The claim is saved with a claim number and the status Submitted, appears in the reviewer's queue, becomes Approved with a history of SUBMIT and APPROVE, and is read back as Approved. The 1,600.00 Meals item is refused under its field (BR-12).}
 {As expected. The claim was \texttt{CLM-2026-000001}, total 650.00 THB; after Approve the toast read ``Claim approved.''; My Claims and the claim detail showed Approved with the two history entries. The Meals item of 1,600.00 gave ``A Meals item cannot be more than 1,500.00 THB.'' under the Amount field, with the input kept.}
 {The row in PostgreSQL after the demos (\texttt{db-readback} below): \texttt{CLM-2026-000001} is \texttt{APPROVED}, version 1, with history \texttt{SUBMIT} then \texttt{APPROVE}. The five Sprint 1 E2E specs (F01-TEST-24, F02-TEST-21, F03-TEST-12, F04-TEST-23, F04-TEST-24) pass on this stack (8 of 8 Playwright specs, Part E).}

\evidencehalf{figures/demo/a1-1-new-claim-two-items.png}{Continuity: \texttt{req.anan} fills New Claim with a Travel item of 250.00 and a Meals item of 400.00, total 650.00 THB (F02-AC-01).}
\evidencehalf{figures/demo/a1-2-claim-submitted.png}{Continuity: the claim is saved with its claim number and the status Submitted (F02-AC-01).}
\evidencehalf{figures/demo/a1-3-reviewer-approves.png}{Continuity: \texttt{rev.malee} approves the claim from her Review Queue (F04-AC-04).}
\evidencehalf{figures/demo/a1-4-my-claims-approved.png}{Continuity: retrieval of the saved result. \texttt{req.anan}'s My Claims shows the claim as Approved (F03-AC-01).}
\evidencehalf{figures/demo/a1-5-claim-detail-approved-with-history.png}{Continuity: the claim detail read back from the database: Approved, items, total and the history entries (F03-AC-05).}
\evidencehalf{figures/demo/a1-6-meals-limit-refused.png}{Continuity: the Sprint 1 rule BR-12 still refuses a Meals item over 1,500.00 THB under its field (F02-AC-06).}

\textbf{Backend read-back of the continuity claim and of the Sprint 2 demonstration claims} (PostgreSQL of the demo stack; the Sprint 2 rows are explained in Part B):

\VerbatimInput[tabsize=4, fontsize=\scriptsize, frame=single, framesep=5pt, rulecolor=\color{black!25}, breaklines=true]{content/data/db-readback.txt}

\partsection{A.2}{Change and defect register}\label{A:2}

Every deployment, UAT and maintenance finding that touched the Sprint 1 baseline, from Labs 5 to 7 and the Sprint 1 release review. IDs are GitHub issue or pull request numbers; dates are UTC as GitHub records them. ``Re-test'' names the executed check. The sources are \path{docs/labs/lab-06/uat-report.md}, \path{docs/labs/lab-07/change-impact.md}, \path{docs/traceability.md} (change log) and the issues and pull requests cited. Lab 6 findings \emph{we raised} on the Queue team's product (Queue \#55 to \#58) belong to that product and are not in this register.

{\footnotesize\setlength{\tabcolsep}{3pt}
\begin{longtable}{@{}p{0.9cm}p{2.6cm}p{2.3cm}p{2.7cm}p{2.0cm}p{1.9cm}p{3.6cm}@{}}
\toprule \textbf{ID} & \textbf{Source} & \textbf{Requirement, feature} & \textbf{Decision} & \textbf{Status} & \textbf{Fix release} & \textbf{Re-test} \\ \midrule \endhead
\#86 & Lab 6 UAT by the Queue team (\texttt{MacOverlorD}), case X-02/X-03, on \texttt{eed7b27}, 2026-11-19 & FR-02, NFR-08, F02-AC-12 (title, purpose, description, receipt reference of only spaces stored blank) & Accepted defect, minor; owner \texttt{beambeambeam} & Resolved; closed 2026-11-22 & \texttt{bb5a0d1} (PR \#92 into \texttt{lab6-staging}, release PR \#93); first tag \texttt{v0.1.0.0} & Red \texttt{f684459}: \texttt{status = 201, want 400}; green \texttt{ecb4b2e}, \texttt{2054923}; retest PASS recorded in the UAT report; F02-TEST-18 and F02-TEST-19 pass in Part B's AC table \\
\#87 & Lab 6 UAT, case X-04, on \texttt{eed7b27}, 2026-11-20 & FR-02, F02-AC-13 (double-click on Submit created CLM-2026-000015 and 000016) & Accepted defect, major; owner \texttt{beambeambeam}; server idempotency left out by design (SDS 7.2) & Resolved; closed 2026-11-22 & \texttt{bb5a0d1} (PR \#91, release PR \#93); \texttt{v0.1.0.0} & Component test red on \texttt{f684459}, green \texttt{f97dec4} to \texttt{496b22f}; retest PASS recorded in the UAT report \\
\#99 & Lab 7 maintenance change L7/01 (PR \#101 contract, PR \#105 code) & FR-06, BR-04, F04-AC-06, F04-AC-18 (reject reason 5 to 500 becomes 10 to 500) & Changed requirement, not a defect; approved by \texttt{Kiatisakk} (PR \#101); OpenAPI review on PR \#105 & Resolved; closed 2026-11-24 & \texttt{c32e645} (release PR \#106); \texttt{v0.1.0.0} & Red \texttt{214d1d5}, green \texttt{b51c058} (PR \#105); revised F04-TEST-06, 07, 21, 24 pass; the same minimum now also guards Return (F06-AC-02, F06-TEST-02), passing in \texttt{v0.1.1.0} \\
\#66 & Sprint 1 release review (PR \#69): E2E specs did not run in CI & Product DoD item 4, NFR-08 evidence & Decision: E2E journeys run on a local seeded stack before each release, not in CI (Product DoD 0.3, PR \#96) & Resolved by decision; closed 2026-11-22 & \texttt{fe9eb01}; \texttt{v0.1.0.0} & 8 of 8 Playwright specs pass on the \texttt{3a91ea4} stack (Part E) \\
\#67, \#70 & Sprint 1 review follow-ups of PR \#65 and release PR \#69 & SDS 3.1 signatures, the 403 access page, session audit fields & Accepted as follow-up work; SDS 0.12 recorded them as as-built notes & Resolved; closed 2026-11-22 & \texttt{fbf4382} (PR \#97); \texttt{v0.1.0.0} & F04-TEST-21 gained the approve 403 case; Go and web suites pass at \texttt{3a91ea4} (Part E) \\
\#61 & Sprint 1 PR \#60 known gap: receipts are a text reference only & EX-05 (receipt upload excluded), F-EXTRA-01 & Not implemented; out of scope for the baseline & Closed 2026-11-24 without an implementing commit (\texttt{git log --grep} finds none); SRS still lists EX-05 & none & none. \todo{kiatisakk: confirm why \#61 was closed and add one sentence} \\
none & Sprint 2 release checks (\texttt{docs/traceability.md}): \texttt{bun run dev} could not start the API on Windows & NFR-10 (README run path) & Fix: air build entry for Windows & Resolved & \texttt{33ade04}; in \texttt{v0.1.1.0} & README steps rerun on a clean clone: air started the API, \texttt{/api/v1/ready} gave 200 \\
\bottomrule
\end{longtable}}

\textbf{Remaining status.} No register item is open. The known limitations that are not findings (a local-only web test failure, no E2E job in CI, no receipt upload) are listed in the back matter. \todo{recapture on candidate: add any finding raised against the candidate after 2026-11-25, and any Lab 7 re-test of the candidate}

\partsection{A.3}{Approved baseline revisions and their rationale}\label{A:3}

The baseline of Deliverable 1 is SRS 0.2, the Sprint 1 contracts at the release \texttt{v0.0.3.0}, SDS 0.12 and Product DoD 0.2. Revisions since, from the revision histories and the change log of \path{docs/traceability.md}. Approver is the name recorded in the document.

{\footnotesize\setlength{\tabcolsep}{3pt}
\begin{longtable}{@{}p{3.0cm}p{1.9cm}p{6.7cm}p{4.6cm}@{}}
\toprule \textbf{Document, version} & \textbf{Date} & \textbf{Change} & \textbf{Rationale, approval} \\ \midrule \endhead
SRS 0.3 & 2026-11-24 & BR-04: the reason minimum rises from 5 to 10 characters for reject, return and final reject; the maximum stays 500 & Lab 7 change L7/01 (\#99): ``No pic'' is accepted today and gives the requester nothing to act on; approved by 1005 \\
F-04 contract 0.6 & 2026-11-22 & Playwright item follows Product DoD 0.3 item 4 (local seeded stack, not CI); no AC or test changed & \#66; \texttt{beambeambeam} \\
F-04 contract 0.7 & 2026-11-22 & F04-TEST-21 adds the approve 403 case & Sprint 1 follow-up \#67; \texttt{beambeambeam} \\
F-04 contract 0.8 & 2026-11-24 & F04-AC-06 and F04-AC-18 and the inputs of F04-TEST-06, 07, 21, 24 revised for BR-04; the original red and green commits apply to the contract before 0.8 & L7/01 (\#99); 1005 \\
F-04 contract 0.9 & 2026-11-24 & Seed gains DRAFT, RETURNED and WITHDRAWN sample claims; F04-TEST-22 expects them & Sprint 2 contracts (PR \#100) \\
Product DoD (Sprint 1) 0.3 & 2026-11-22 & Item 4: E2E specs pass on a local seeded stack before each release, tested commit recorded & The course specifications require E2E journeys to exist, not to run in CI (\#66); \texttt{beambeambeam} \\
System SDS 0.13 to 0.15 & 2026-11-17 to 11-19 & Lab 5: root Dockerfile and \texttt{WEB\_DIR}; \texttt{image} required check; rulesets; release identity is the SHA & Lab 5 (\#53, \#54); 1005 \\
System SDS 0.16, 0.17 & 2026-11-22 & E2E strategy (10); session audit fields and \texttt{Create} as built & \#66, \#67, \#70; \texttt{beambeambeam} \\
System SDS 0.18 & 2026-11-24 & BR-04 minimum in \texttt{claim\_history.comment} is 10; no schema change & L7/01; 1005 \\
Sprint plan (change log) & 2026-11-15, 2026-11-23 & Sprint 1 finished 15 Nov, not 8 Nov; Sprint 2 started 23 Nov from \texttt{sprint-2} & Schedule shift recorded in the plan; the log marks the first two staging-branch entries ``Pending review'' \\
Traceability change log & 2026-11-13 to 11-25 & One dated row per change of mapping or recorded result, including the Sprint 2 AC and tests (2026-11-23) and the Sprint 2 release checks (2026-11-25) & Each row names its ticket or PR; shown in full in Part D \\
\bottomrule
\end{longtable}}

\noindent The Sprint 2 contracts F-05, F-06 and F-07 have their own revision histories (Part B, 0.1 to 0.4); those are Sprint 2 work, not changes to the Sprint 1 baseline.

\reportpartend{A}

\reportpart{B}{Sprint 2 contracts and working increment}
% Generated by scripts/gen-part-lists.ts. Do not edit.

\partsection{F-05 contract}{}
\docsource{docs/sprints/sprint-2/F-05-drafts.md}{main}{3a91ea4}
\evidencepart{figures/docs/f05-p1.png}{\detokenize{docs/sprints/sprint-2/F-05-drafts.md, piece 1 of 12}}
\evidencepart{figures/docs/f05-p2.png}{\detokenize{docs/sprints/sprint-2/F-05-drafts.md, piece 2 of 12}}
\evidencepart{figures/docs/f05-p3.png}{\detokenize{docs/sprints/sprint-2/F-05-drafts.md, piece 3 of 12}}
\evidencepart{figures/docs/f05-p4.png}{\detokenize{docs/sprints/sprint-2/F-05-drafts.md, piece 4 of 12}}
\evidencepart{figures/docs/f05-p5.png}{\detokenize{docs/sprints/sprint-2/F-05-drafts.md, piece 5 of 12}}
\evidencepart{figures/docs/f05-p6.png}{\detokenize{docs/sprints/sprint-2/F-05-drafts.md, piece 6 of 12}}
\evidencepart{figures/docs/f05-p7.png}{\detokenize{docs/sprints/sprint-2/F-05-drafts.md, piece 7 of 12}}
\evidencepart{figures/docs/f05-p8.png}{\detokenize{docs/sprints/sprint-2/F-05-drafts.md, piece 8 of 12}}
\evidencepart{figures/docs/f05-p9.png}{\detokenize{docs/sprints/sprint-2/F-05-drafts.md, piece 9 of 12}}
\evidencepart{figures/docs/f05-p10.png}{\detokenize{docs/sprints/sprint-2/F-05-drafts.md, piece 10 of 12}}
\evidencepart{figures/docs/f05-p11.png}{\detokenize{docs/sprints/sprint-2/F-05-drafts.md, piece 11 of 12}}
\evidencepart{figures/docs/f05-p12.png}{\detokenize{docs/sprints/sprint-2/F-05-drafts.md, piece 12 of 12}}
\noindent\textit{The 2 diagrams of \path{docs/sprints/sprint-2/F-05-drafts.md}, repeated at full size from the same commit (3a91ea4):}
\diagramfull{figures/diagrams/f05-1.png}{\detokenize{docs/sprints/sprint-2/F-05-drafts.md}, diagram 1 of 2}
\diagramfull{figures/diagrams/f05-2.png}{\detokenize{docs/sprints/sprint-2/F-05-drafts.md}, diagram 2 of 2}

\partsection{F-06 contract}{}
\docsource{docs/sprints/sprint-2/F-06-return-resubmit.md}{main}{3a91ea4}
\evidencepart{figures/docs/f06-p1.png}{\detokenize{docs/sprints/sprint-2/F-06-return-resubmit.md, piece 1 of 9}}
\evidencepart{figures/docs/f06-p2.png}{\detokenize{docs/sprints/sprint-2/F-06-return-resubmit.md, piece 2 of 9}}
\evidencepart{figures/docs/f06-p3.png}{\detokenize{docs/sprints/sprint-2/F-06-return-resubmit.md, piece 3 of 9}}
\evidencepart{figures/docs/f06-p4.png}{\detokenize{docs/sprints/sprint-2/F-06-return-resubmit.md, piece 4 of 9}}
\evidencepart{figures/docs/f06-p5.png}{\detokenize{docs/sprints/sprint-2/F-06-return-resubmit.md, piece 5 of 9}}
\evidencepart{figures/docs/f06-p6.png}{\detokenize{docs/sprints/sprint-2/F-06-return-resubmit.md, piece 6 of 9}}
\evidencepart{figures/docs/f06-p7.png}{\detokenize{docs/sprints/sprint-2/F-06-return-resubmit.md, piece 7 of 9}}
\evidencepart{figures/docs/f06-p8.png}{\detokenize{docs/sprints/sprint-2/F-06-return-resubmit.md, piece 8 of 9}}
\evidencepart{figures/docs/f06-p9.png}{\detokenize{docs/sprints/sprint-2/F-06-return-resubmit.md, piece 9 of 9}}
\noindent\textit{The 1 diagram of \path{docs/sprints/sprint-2/F-06-return-resubmit.md}, repeated at full size from the same commit (3a91ea4):}
\diagramfull{figures/diagrams/f06-1.png}{\detokenize{docs/sprints/sprint-2/F-06-return-resubmit.md}, diagram 1 of 1}

\partsection{F-07 contract}{}
\docsource{docs/sprints/sprint-2/F-07-withdraw.md}{main}{3a91ea4}
\evidencepart{figures/docs/f07-p1.png}{\detokenize{docs/sprints/sprint-2/F-07-withdraw.md, piece 1 of 6}}
\evidencepart{figures/docs/f07-p2.png}{\detokenize{docs/sprints/sprint-2/F-07-withdraw.md, piece 2 of 6}}
\evidencepart{figures/docs/f07-p3.png}{\detokenize{docs/sprints/sprint-2/F-07-withdraw.md, piece 3 of 6}}
\evidencepart{figures/docs/f07-p4.png}{\detokenize{docs/sprints/sprint-2/F-07-withdraw.md, piece 4 of 6}}
\evidencepart{figures/docs/f07-p5.png}{\detokenize{docs/sprints/sprint-2/F-07-withdraw.md, piece 5 of 6}}
\evidencepart{figures/docs/f07-p6.png}{\detokenize{docs/sprints/sprint-2/F-07-withdraw.md, piece 6 of 6}}
\noindent\textit{The 1 diagram of \path{docs/sprints/sprint-2/F-07-withdraw.md}, repeated at full size from the same commit (3a91ea4):}
\diagramfull{figures/diagrams/f07-1.png}{\detokenize{docs/sprints/sprint-2/F-07-withdraw.md}, diagram 1 of 1}

\partsection{Sprint 2 Definition of Done}{}
\docsource{docs/sprints/sprint-2/definition-of-done.md}{main}{3a91ea4}
\evidencepart{figures/docs/dod-p1.png}{\detokenize{docs/sprints/sprint-2/definition-of-done.md, piece 1 of 2}}
\evidencepart{figures/docs/dod-p2.png}{\detokenize{docs/sprints/sprint-2/definition-of-done.md, piece 2 of 2}}

\partsection{B.2}{Sprint 2 completion}\label{B:2}

All three Sprint 2 features were implemented, reviewed and integrated into \texttt{sprint-2} and then released to \texttt{main} as \texttt{v0.1.1.0} (\texttt{3a91ea4}, release PR \#109). No Sprint 2 feature or AC moved to Sprint 3. Contracts were written first: PR \#100 (\texttt{dfc5f6f}, 2026-11-24 22:32 UTC+7) merged before the first implementation commit of PR \#104 (\texttt{fad37c0}, 2026-11-24 22:36 UTC+7, a red test).

\begin{center}\small\setlength{\tabcolsep}{3pt}
\begin{tabular}{@{}p{2.6cm}p{1.9cm}p{2.4cm}p{2.8cm}p{2.1cm}p{3.1cm}@{}}
\toprule \textbf{Feature} & \textbf{Contract} & \textbf{AC} & \textbf{Planned tests} & \textbf{Implemented in} & \textbf{Tickets and release} \\ \midrule
F-05 Drafts & revision 0.3 & 20 (F05-AC-01 to 20) & 26 (F05-TEST-01 to 26) & PR \#104 & \#39 to \#42; \texttt{v0.1.1.0} \\
F-06 Return and resubmit & revision 0.4 & 13 (F06-AC-01 to 13) & 17 (F06-TEST-01 to 17) & PR \#108 & \#43, \#44; \texttt{v0.1.1.0} \\
F-07 Withdraw & revision 0.2 & 10 (F07-AC-01 to 10) & 13 (F07-TEST-01 to 13) & PR \#108 & \#45; \texttt{v0.1.1.0} \\
\bottomrule
\end{tabular}
\end{center}

\textbf{Completion table.} One row per acceptance criterion of the three contracts (the criteria themselves are in B.1). \emph{Covers} is the requirement column of the contract's AC table. \emph{Implementation} names the code unit under \path{apps/api/internal/claim}, the endpoint, or the web component that delivers the AC. \emph{Executed verification} lists the planned tests of the contract's Software Test Specification that cover the AC, with their level. \emph{Result} counts those tests that passed in this report's runs (Go: the run on \texttt{3a91ea4} in B.4; Vitest: CI run 45614889848 on \texttt{3a91ea4}, 182 of 182 passed; Playwright: local seeded stack at \texttt{3a91ea4}, 8 of 8 passed). \emph{Accepted in} is the release that carries the AC.
\texttt{F0N-T0M} abbreviates \texttt{F0N-TEST-0M}. The table is generated by \path{docs/reports/scripts/gen-ac-table.ts} from the contracts, the test sources and \path{go.json}.

% Generated by scripts/gen-ac-table.ts at 3a91ea4. Do not edit.
{\footnotesize\setlength{\tabcolsep}{3pt}
\begin{longtable}{@{}p{1.6cm}p{1.9cm}p{4.6cm}p{3.9cm}p{2.9cm}p{1.6cm}@{}}
\toprule \textbf{AC} & \textbf{Covers} & \textbf{Implementation} & \textbf{Executed verification} & \textbf{Result} & \textbf{Accepted in} \\ \midrule \endhead
F05-AC-01 & FR-08, FR-03, BR-08 & POST /claims?draft=true: Handler.CreateClaim, Service.Create with submit=false (internal/claim) & F05-T01 (Integration) & pass (1 of 1) & v0.1.1.0 \\
F05-AC-02 & BR-05, BR-06, BR-12, NFR-08 & rules.go: item and total limits for drafts, dates skipped until submit & F05-T02 (Integration (module)) & pass (1 of 1) & v0.1.1.0 \\
F05-AC-03 & BR-15 & policy.go and Service.Create: BR-15 guard, no reviewer means ErrForbidden & F05-T03, F05-T04 (Integration, Integration (module)) & pass (2 of 2) & v0.1.1.0 \\
F05-AC-04 & FR-08, NFR-02 & policy.go: owner-only access to a DRAFT; allowedActions from Allowed & F05-T05, F05-T06 (Integration, Unit) & pass (2 of 2) & v0.1.1.0 \\
F05-AC-05 & FR-08, FR-03, BR-05, BR-06, BR-08, BR-12 & PUT /claims/\{id\}: Handler.UpdateClaim, edit.go (replace items, recompute total) & F05-T07, F05-T08 (Integration, Integration (module)) & pass (2 of 2) & v0.1.1.0 \\
F05-AC-06 & BR-10, NFR-02 & policy.go owner check; lifecycle.go BR-10 on a non-draft & F05-T06, F05-T09 (Unit, Integration) & pass (2 of 2) & v0.1.1.0 \\
F05-AC-07 & BR-13 & Handler If-Match parsing (428, 400); version compare in Service & F05-T10 (Integration) & pass (1 of 1) & v0.1.1.0 \\
F05-AC-08 & FR-08, FR-07, BR-17, NFR-06 & POST /claims/\{id\}/submit: Handler.SubmitClaim, Service DRAFT to SUBMITTED with SUBMIT history & F05-T11 (Integration) & pass (1 of 1) & v0.1.1.0 \\
F05-AC-09 & BR-07 & rules.go: expense-date window counted from the submission time & F05-T12, F05-T13 (Integration (module), Integration) & pass (2 of 2) & v0.1.1.0 \\
F05-AC-10 & BR-17 & Service submit reads the owner's current reviewer (BR-17) & F05-T14 (Integration (module)) & pass (1 of 1) & v0.1.1.0 \\
F05-AC-11 & BR-15 & policy.go: no submit without a reviewer; Service submit guard (BR-15) & F05-T06, F05-T15 (Unit, Integration (module)) & pass (2 of 2) & v0.1.1.0 \\
F05-AC-12 & BR-03, BR-13 & lifecycle.go transitions (BR-03); conditional UPDATE on version & F05-T16, F05-T17 (Integration, Integration (module)) & pass (2 of 2) & v0.1.1.0 \\
F05-AC-13 & NFR-06 & Service: status update and history insert in one transaction & F05-T18 (Integration (module)) & pass (1 of 1) & v0.1.1.0 \\
F05-AC-14 & FR-08, BR-10 & DELETE /claims/\{id\}: Handler.DeleteClaim, Service delete for a DRAFT of the owner & F05-T19 (Integration) & pass (1 of 1) & v0.1.1.0 \\
F05-AC-15 & BR-10, BR-11 & migration 00002\_claims.sql: history foreign key and append-only trigger & F05-T20 (Integration) & pass (1 of 1) & v0.1.1.0 \\
F05-AC-16 & FR-08, NFR-08 & web: new-claim-page.tsx, claim-form.tsx (Save as draft) & F05-T21 (Component) & pass (1 of 1) & v0.1.1.0 \\
F05-AC-17 & FR-08, NFR-08 & web: claim-detail-page.tsx, claim-owner-actions.tsx & F05-T22 (Component) & pass (1 of 1) & v0.1.1.0 \\
F05-AC-18 & FR-08, NFR-08 & web: edit-claim-page.tsx and the claims/\$claimId\_.edit route & F05-T23 (Component) & pass (1 of 1) & v0.1.1.0 \\
F05-AC-19 & FR-08 & web: my-claims-page.tsx, status-badge.tsx & F05-T24 (Component) & pass (1 of 1) & v0.1.1.0 \\
F05-AC-20 & FR-08, FR-07 & web/e2e/f05-drafts.spec.ts over the API and web above & F05-T25 (E2E) & pass (1 of 1) & v0.1.1.0 \\
F06-AC-01 & FR-09, FR-07, NFR-06 & POST /claims/\{id\}/return: Handler.ReturnClaim, Service.Act with ActionReturn & F06-T01 (Integration) & pass (1 of 1) & v0.1.1.0 \\
F06-AC-02 & BR-04 & rules.go: reason length shared with Reject (MinReasonLength, MaxCommentLength) & F06-T02, F06-T03 (Integration (module), Integration) & pass (2 of 2) & v0.1.1.0 \\
F06-AC-03 & BR-01, BR-02 & policy.go: assigned-reviewer relation (BR-01, BR-02, ADR 0001) & F06-T04 (Integration) & pass (1 of 1) & v0.1.1.0 \\
F06-AC-04 & BR-03, BR-13 & lifecycle.go transitions (BR-03); version compare (BR-13) & F06-T05 (Integration) & pass (1 of 1) & v0.1.1.0 \\
F06-AC-05 & BR-14 & policy.go: MaxReturns = 3 hides return (BR-14); Service.Act refuses it & F06-T06, F06-T07 (Unit, Integration (module)) & pass (2 of 2) & v0.1.1.0 \\
F06-AC-06 & FR-09, FR-10 & list scopes in Service.List; policy for the reviewer view & F06-T08 (Integration) & pass (1 of 1) & v0.1.1.0 \\
F06-AC-07 & FR-10, FR-07, BR-17, NFR-06 & Service submit from RETURNED keeps the reviewer and firstSubmittedAt & F06-T09, F06-T17 (Integration, Integration (module)) & pass (2 of 2) & v0.1.1.0 \\
F06-AC-08 & BR-07 & rules.go: 90-day window counted from firstSubmittedAt (BR-07) & F06-T10 (Integration (module)) & pass (1 of 1) & v0.1.1.0 \\
F06-AC-09 & FR-10, FR-03, BR-10 & edit.go and policy.go: owner edits a RETURNED claim & F06-T11 (Integration) & pass (1 of 1) & v0.1.1.0 \\
F06-AC-10 & NFR-06 & Service: status update and history insert in one transaction & F06-T12 (Integration (module)) & pass (1 of 1) & v0.1.1.0 \\
F06-AC-11 & FR-09, NFR-08 & web: decision-panel.tsx (Return dialog, Returned n / 3) & F06-T13 (Component) & pass (1 of 1) & v0.1.1.0 \\
F06-AC-12 & FR-10, NFR-08 & web: claim-detail-page.tsx banner, claim-owner-actions.tsx Resubmit & F06-T14 (Component) & pass (1 of 1) & v0.1.1.0 \\
F06-AC-13 & FR-09, FR-10, FR-07 & web/e2e/f06-return-resubmit.spec.ts over the API and web above & F06-T15 (E2E) & pass (1 of 1) & v0.1.1.0 \\
F07-AC-01 & FR-11, FR-07, NFR-06 & POST /claims/\{id\}/withdraw: Handler.WithdrawClaim, Service.Act with ActionWithdraw & F07-T01 (Integration) & pass (1 of 1) & v0.1.1.0 \\
F07-AC-02 & FR-11 & openapi.yaml ClaimWithdrawal.comment (x-max-length 500); validate.go & F07-T02 (Integration) & pass (1 of 1) & v0.1.1.0 \\
F07-AC-03 & BR-16 & lifecycle.go transitions: only SUBMITTED and RETURNED reach WITHDRAWN (BR-16) & F07-T03 (Integration) & pass (1 of 1) & v0.1.1.0 \\
F07-AC-04 & FR-11, NFR-02 & policy.go: ownership check for withdraw & F07-T04, F07-T05 (Integration, Unit) & pass (2 of 2) & v0.1.1.0 \\
F07-AC-05 & BR-13 & Handler If-Match parsing (428, 400); version compare in Service & F07-T06 (Integration) & pass (1 of 1) & v0.1.1.0 \\
F07-AC-06 & BR-03, BR-13 & lifecycle.go: WITHDRAWN has no outgoing transition; policy gives [] & F07-T05, F07-T07, F07-T08 (Unit, Integration, Integration (module)) & pass (3 of 3) & v0.1.1.0 \\
F07-AC-07 & NFR-06 & Service: status update and history insert in one transaction & F07-T09 (Integration (module)) & pass (1 of 1) & v0.1.1.0 \\
F07-AC-08 & FR-11 & Service.List review scope excludes WITHDRAWN; Get still allowed & F07-T10 (Integration) & pass (1 of 1) & v0.1.1.0 \\
F07-AC-09 & FR-11, NFR-08 & web: claim-owner-actions.tsx (Withdraw dialog), status-badge.tsx & F07-T11 (Component) & pass (1 of 1) & v0.1.1.0 \\
F07-AC-10 & FR-11 & web/e2e/f07-withdraw.spec.ts over the API and web above & F07-T12 (E2E) & pass (1 of 1) & v0.1.1.0 \\
\bottomrule\end{longtable}}


\noindent Totals: 43 of 43 AC have at least one executed passing test; 56 of 56 planned tests have an actual test and pass (0 failed, 0 skipped, 0 manual). The seed rows F05-TEST-26, F06-TEST-16 and F07-TEST-13 are subtests of \texttt{TestF04TEST22SeedAddsSampleClaimsOnce} (they verify F04-AC-19, the seed, not a Sprint 2 AC).

\partsection{B.3}{Program demonstrations}\label{B:3}

Three scenarios on the Sprint 2 release, run on the local stack described in A.1 (tag \texttt{v0.1.1.0}, commit \texttt{3a91ea4}, 2026-11-25, Chrome, synthetic data), taken by \path{docs/reports/sprints/02/demo.ts}, not by hand.
Between them they show saved data read back from PostgreSQL (Demo 1, A.1), meaningful rule rejections (BR-07 at draft submission, BR-04 on the return reason, BR-14 after three returns) and backend evidence of a forbidden operation (a 403 from the API).
The candidate is \texttt{v0.1.1.0} itself, so these demonstrations were run on the candidate; no later code exists on \texttt{main} (E.1).

\subsection*{Demo 1: F-05 Drafts: save, read back, rule rejection, ownership 403}

\demofacts
 {F-05; F05-AC-01, F05-AC-04, F05-AC-08, F05-AC-09, F05-AC-16 to F05-AC-19; BR-07, NFR-02; F05-TEST-01, 05, 11, 12, 13}
 {Requester \texttt{req.anan}; then requester \texttt{req.ploy} for the 403}
 {Fresh seeded database; \texttt{req.anan} has a reviewer (\texttt{rev.malee})}
 {(1) \texttt{req.anan} fills New Claim with one Travel item and presses \textbf{Save as draft}. (2) Opens My Claims, then the API URL of the draft. (3) Edits the draft and sets the expense date to tomorrow; Save changes; presses \textbf{Submit claim}. (4) Sets the date to yesterday; Submit claim. (5) Saves a second draft; logs in as \texttt{req.ploy} and opens that draft's API URL.}
 {(1) The API creates the draft (201) with status Draft and no history. (2) My Claims lists it with the Draft badge; the API returns it with \texttt{status DRAFT}, \texttt{version 0}, \texttt{assignedReviewer null}. (3) Saving with a future date is accepted (drafts skip dates); submitting is refused with 422 ``Line 1: Expense date cannot be in the future.'' and the claim stays a draft. (4) The corrected draft is submitted: status Submitted, history \texttt{Draft} to \texttt{Submitted}. (5) The other requester gets 403 with no claim data.}
 {As expected. Draft \texttt{CLM-2026-000002} (250.00 THB) was saved with the toast ``Draft saved.'', listed in My Claims, and returned by \texttt{GET /api/v1/claims/\{id\}} with \texttt{"status":"DRAFT"}, \texttt{"version":0}, \texttt{"history":[]}. The submit with a date of 26 Nov 2026 was refused with ``Check the claim before submitting. Line 1: Expense date cannot be in the future.''; after the correction the claim became Submitted (version 3). \texttt{req.ploy}'s read of \texttt{CLM-2026-000003} returned \texttt{\{"code":"FORBIDDEN","status":403,\ldots\}}.}
 {Saved data read back from PostgreSQL: the API answer in the third figure, and the history row \texttt{SUBMIT DRAFT to SUBMITTED} in the read-back of A.1. The 403 is the API's own response, not a hidden button: automated F05-TEST-05 (\texttt{TestF05TEST05OnlyOwnerCanOpenDraft}) and F05-TEST-03 (Finance Officer 403) pass (table below).}

\evidencehalf{figures/demo/d1-1-draft-saved.png}{Demo 1: the draft is saved with the Draft badge, the server total and no history (F05-AC-01, F05-AC-16).}
\evidencehalf{figures/demo/d1-2-my-claims-lists-the-draft.png}{Demo 1: My Claims lists the draft (F05-AC-19).}
\evidencesmall{figures/demo/d1-3-api-returns-the-saved-draft.png}{Demo 1: retrieval of saved data. The API's own JSON for the draft, read from PostgreSQL: \texttt{status DRAFT}, \texttt{version 0}, \texttt{history []}, the owner's \texttt{allowedActions} edit, delete, submit (F05-AC-01, F05-AC-04).}{0.96}
\evidencehalf{figures/demo/d1-4-submit-refused-future-date.png}{Demo 1: rule rejection. Submitting the draft with a future expense date is refused under BR-07 and the claim stays a draft (F05-AC-09, F05-AC-17).}
\evidencehalf{figures/demo/d1-5-draft-submitted.png}{Demo 1: the corrected draft is submitted: Submitted, one history entry from Draft (F05-AC-08).}
\evidencesmall{figures/demo/d1-6-other-requester-gets-403.png}{Demo 1: backend evidence of a forbidden ownership operation. \texttt{req.ploy} reads \texttt{req.anan}'s draft through the API: 403 \texttt{FORBIDDEN} with no claim data (F05-AC-04, NFR-02).}{0.96}

\subsection*{Demo 2: F-06 Return and resubmit: reason rule, return, resubmit, return limit}

\demofacts
 {F-06; F06-AC-01, F06-AC-02, F06-AC-05, F06-AC-07, F06-AC-11, F06-AC-12; BR-04, BR-14; F06-TEST-02, 07, 13, 14, 15}
 {Requester \texttt{req.anan}; reviewer \texttt{rev.malee}}
 {Fresh seeded database; \texttt{req.anan}'s Submitted claim in \texttt{rev.malee}'s queue}
 {(1) \texttt{rev.malee} opens the claim, chooses Return and types ``Too short''. (2) Types ``Please attach the taxi receipt'' and confirms. (3) \texttt{req.anan} opens the claim, reads the reason and ``Returned 1 / 3'', and presses Resubmit. (4) Two more return and resubmit rounds. (5) \texttt{rev.malee} opens the claim a fourth time.}
 {(1) The dialog refuses a reason under 10 characters (BR-04): message and a disabled Return button. (2) The claim becomes Returned with the reason in its history. (3) The owner sees the banner, the reason and the count; Resubmit makes it Submitted again for the same reviewer. (5) After three returns Return is no longer offered; Approve and Reject stay (BR-14).}
 {As expected. ``Reason must be at least 10 characters.'' appeared at 9/500 characters with Return disabled. The claim \texttt{CLM-2026-000004} was returned, resubmitted, and after round 3 the reviewer's panel read ``Returned 3 / 3 -- approve or reject.'' with only Approve and Reject.}
 {The database after the run: \texttt{CLM-2026-000004} is \texttt{SUBMITTED}, \texttt{version 6}, \texttt{return\_count 3}, with history SUBMIT, then RETURN and RESUBMIT three times (read-back in A.1; the return comments are stored). The refusal of a fourth return by the API (409 \texttt{INVALID\_TRANSITION}, BR-14) is F06-TEST-07, shown with its output in B.4. The CHECK \texttt{return\_count <= 3} is in the schema.}

\evidencehalf{figures/demo/d2-1-short-reason-refused.png}{Demo 2: rule rejection. A 9-character return reason is refused under BR-04 and the Return button stays disabled (F06-AC-02, F06-AC-11).}
\evidencehalf{figures/demo/d2-2-claim-returned.png}{Demo 2: the reviewer returns the claim with a valid reason (F06-AC-01).}
\evidencehalf{figures/demo/d2-3-owner-sees-reason-and-count.png}{Demo 2: the owner sees the return reason, ``Returned 1 / 3'' and the Edit and Resubmit actions (F06-AC-12).}
\evidencehalf{figures/demo/d2-4-claim-resubmitted.png}{Demo 2: Resubmit puts the claim back to the same reviewer (F06-AC-07).}
\evidencehalf{figures/demo/d2-5-no-fourth-return.png}{Demo 2: after three returns the reviewer has no Return action, only Approve and Reject; the history shows Return and Resubmit rounds (F06-AC-05, BR-14).}

\subsection*{Demo 3: F-07 Withdraw}

\demofacts
 {F-07; F07-AC-01, F07-AC-08, F07-AC-09; BR-16; F07-TEST-01, 10, 11, 12}
 {Requester \texttt{req.anan}; reviewer \texttt{rev.malee}}
 {A Submitted claim of \texttt{req.anan}}
 {\texttt{req.anan} opens the claim, chooses Withdraw, types the comment ``No longer needed'' and confirms; \texttt{rev.malee} opens her Review Queue.}
 {The claim becomes Withdrawn with the comment in its history and no actions; it is gone from the reviewer's queue.}
 {As expected: ``Claim withdrawn.'', the Withdrawn badge and the history entry \texttt{Submitted -> Withdrawn} with the comment; the Review Queue table did not list the claim.}
 {\texttt{CLM-2026-000005} is \texttt{WITHDRAWN}, \texttt{version 1}, with history SUBMIT and WITHDRAW (``No longer needed'') in the A.1 read-back. Access: only the owner may withdraw (F07-TEST-04, B.4); withdrawing from a non-submitted, non-returned status is refused with 409 BR-16 (F07-TEST-03, B.4).}

\evidencehalf{figures/demo/d3-1-withdraw-dialog.png}{Demo 3: the Withdraw dialog with the optional comment and its counter (F07-AC-09).}
\evidencehalf{figures/demo/d3-2-claim-withdrawn-with-history.png}{Demo 3: the claim is Withdrawn; the history shows the withdrawal and its comment (F07-AC-01).}
\evidencehalf{figures/demo/d3-3-review-queue-without-the-claim.png}{Demo 3: the claim is no longer in \texttt{rev.malee}'s Review Queue (F07-AC-08).}

\textbf{Backend evidence for the access restrictions of Sprint 2.} Go suite on \texttt{3a91ea4} (B.4); every test passes. Each sends the request through the real router and checks the stored state is unchanged.

\begin{center}\small
\begin{tabular}{@{}llp{8.4cm}l@{}}
\toprule \textbf{Test} & \textbf{AC} & \textbf{Asserts} & \textbf{Result} \\ \midrule
F05-TEST-03 & F05-AC-03 & \texttt{fin.wanida} saves a draft: 403; no claim stored (BR-15) & pass \\
F05-TEST-05 & F05-AC-04 & only the owner opens a draft: \texttt{req.ploy}, the owner's reviewer and the Finance Officer get 403 with no claim data & pass \\
F05-TEST-09 & F05-AC-06 & a non-owner edit gets 403; the owner's edit of a Submitted claim gets 409 BR-10 & pass \\
F06-TEST-04 & F06-AC-03 & only the assigned reviewer returns: the owner and another reviewer get 403 & pass \\
F07-TEST-04 & F07-AC-04 & only the owner withdraws: the assigned reviewer and another requester get 403, no claim data in the body & pass \\
\bottomrule
\end{tabular}
\end{center}

\partsection{B.4}{Verification: commands, results, AC to test, excerpts, TDD}\label{B:4}

\subsubsection*{Commands, tested commit and counts}

Every local command below ran on 2026-11-25 on the tree of \texttt{v0.1.1.0} (\texttt{3a91ea4}, a clean worktree at that commit) on the author's Windows machine with Docker Desktop (PostgreSQL 17 containers through \path{internal/pgtest}). CI ran the same suites on \texttt{3a91ea4} itself: CI run 45614889848 (push to \texttt{main}), every job \texttt{success} (\texttt{image}, \texttt{code\_quality}, \texttt{test}, \texttt{build}, \texttt{api}, \texttt{release}), followed by Deploy demo run 45615736377 (Part E).
\todo{recapture on candidate: rerun the commands on the candidate commit and replace the counts and run IDs}

\textbf{Go API: tests.} \texttt{bun run go-test 02} runs \texttt{go test -json -count=1 ./...} in \path{apps/api} and keeps the JSON stream that feeds the AC table.

\begin{shell}
$ cd apps/api && go test -json -count=1 ./...        (via: bun run go-test 02)
packages with tests: 5   (apps/api, internal/auth, internal/claim, internal/money, migrations)
top-level tests: 178 passed, 0 failed, 0 skipped
subtests:        241 passed, 0 failed, 0 skipped
per package (top-level tests passed): apps/api 114, internal/claim 47, internal/auth 11,
                                      internal/money 3, migrations 3
exit status 0   (6 packages have no test files: env, internal/api, internal/db,
                 internal/httperr, internal/pgtest, internal/system)
\end{shell}

\textbf{Go API: race detector and the 100\% coverage gate.}

\begin{shell}
$ cd apps/api && bun run test:coverage
  = go test -race -coverpkg=./... -coverprofile=coverage.out ./... && go-test-coverage --config=.testcoverage.yml
ok  apps/api                       196.141s
ok  apps/api/internal/auth           2.267s
ok  apps/api/internal/claim        110.250s
ok  apps/api/internal/money          5.935s
ok  apps/api/migrations             35.805s
File coverage threshold (100%) satisfied:     PASS
Package coverage threshold (100%) satisfied:  PASS
Total coverage threshold (100%) satisfied:    PASS
Total test coverage: 100% (956/956)           exit status 0, no data race reported
\end{shell}

\textbf{Web (Vitest), from CI.} The full web suite and its coverage are not run on the author's machine: a full local run times out under load, so CI is the verdict (CI run 45614889848, job \texttt{test}, \texttt{bun run test:coverage}).

\begin{shell}
CI run 45614889848 on 3a91ea4, job test:
 Test Files  22 passed (22)
      Tests  182 passed (182)
 Statements : 100% ( 537/537 )    Branches : 100% ( 313/313 )
 Functions  : 100% ( 182/182 )    Lines    : 100% ( 528/528 )
CI job api: golangci-lint 0 issues; steps Check generated code, Check modules are tidy and
            Check vulnerabilities passed;
            Total test coverage: 100% (956/956)
\end{shell}

\textbf{End-to-end (Playwright, local seeded stack, Chrome).} The specs run on a local stack before each release, not in CI (Product DoD 0.3 item 4). Playwright's own Chromium was not installed on this machine, so the run used the installed Google Chrome through a temporary Playwright config that was not committed; the specs are unchanged.

\begin{shell}
$ SEED_USER_PASSWORD=<set in the process environment> playwright test --timeout=60000   (apps/web/e2e)
  F01-TEST-24, F02-TEST-21, F03-TEST-12, F04-TEST-23, F04-TEST-24,
  F05-TEST-25, F06-TEST-15, F07-TEST-12
  8 passed (1.4m)
\end{shell}

\textbf{Failed and skipped.} Go: 0 failed, 0 skipped (178 top-level tests, 241 subtests). Web in CI: 182 of 182. E2E: 8 of 8. No test was disabled, skipped or weakened to obtain these results. The one local-only web failure recorded in the contracts (\texttt{F01-TEST-20} in \path{auth-routing.test.tsx}, failing alone on the author's machine and passing in CI) is listed in the back matter.

% Generated by scripts/gen-ac-table.ts at 3a91ea4. Do not edit.

\subsubsection*{Automated tests (56 planned tests with an actual test)}
{\footnotesize\setlength{\tabcolsep}{3pt}
\begin{longtable}{@{}p{1.9cm}p{2.3cm}p{1.8cm}p{6.6cm}p{3.6cm}@{}}
\toprule \textbf{Planned test} & \textbf{AC} & \textbf{Level} & \textbf{Actual test (file, name)} & \textbf{Result} \\ \midrule \endhead
F05-TEST-01 & F05-AC-01 & Integration & \path{api/claims_drafts_test.go}\newline \hspace*{1em}\seqsplit{TestF05TEST01SaveDraftWithoutSubmission} & pass (1 run) \\
F05-TEST-02 & F05-AC-02 & Integration (module) & \path{api/internal/claim/drafts_test.go}\newline \hspace*{1em}\seqsplit{TestF05TEST02DraftsCheckItemAndTotalLimitsButSkipDates} & pass (1 run) \\
F05-TEST-03 & F05-AC-03 & Integration & \path{api/claims_drafts_test.go}\newline \hspace*{1em}\seqsplit{TestF05TEST03FinanceOfficerCannotSaveDraft} & pass (1 run) \\
F05-TEST-04 & F05-AC-03 & Integration (module) & \path{api/internal/claim/drafts_test.go}\newline \hspace*{1em}\seqsplit{TestF05TEST04CreateDraftRequiresCurrentReviewer} & pass (1 run) \\
F05-TEST-05 & F05-AC-04 & Integration & \path{api/claims_drafts_test.go}\newline \hspace*{1em}\seqsplit{TestF05TEST05OnlyOwnerCanOpenDraft} & pass (1 run) \\
F05-TEST-06 & F05-AC-04, F05-AC-06, F05-AC-11 & Unit & \path{api/internal/claim/policy_test.go}\newline \hspace*{1em}\seqsplit{TestF05TEST06DraftAndReturnedPolicy} & pass (1 run) \\
F05-TEST-07 & F05-AC-05 & Integration & \path{api/claims_drafts_test.go}\newline \hspace*{1em}\seqsplit{TestF05TEST07EditReplacesItemsAndRecalculatesTotal} & pass (1 run) \\
F05-TEST-08 & F05-AC-05 & Integration (module) & \path{api/internal/claim/drafts_test.go}\newline \hspace*{1em}\seqsplit{TestF05TEST08InvalidEditLeavesDraftUnchanged} & pass (1 run) \\
F05-TEST-09 & F05-AC-06 & Integration & \path{api/claims_drafts_test.go}\newline \hspace*{1em}\seqsplit{TestF05TEST09EditRequiresOwnerAndEditableStatus} & pass (1 run) \\
F05-TEST-10 & F05-AC-07 & Integration & \path{api/claims_drafts_test.go}\newline \hspace*{1em}\seqsplit{TestF05TEST10DraftMutationsRequireCurrentIfMatch} & pass (1 run) \\
F05-TEST-11 & F05-AC-08 & Integration & \path{api/claims_drafts_test.go}\newline \hspace*{1em}\seqsplit{TestF05TEST11SubmitDraftRecordsFirstSubmissionAndQueuesClaim} & pass (1 run) \\
F05-TEST-12 & F05-AC-09 & Integration (module) & \path{api/internal/claim/drafts_test.go}\newline \hspace*{1em}\seqsplit{TestF05TEST12SubmissionChecksExpenseDateWindow} & pass (1 run) \\
F05-TEST-13 & F05-AC-09 & Integration & \path{api/claims_drafts_test.go}\newline \hspace*{1em}\seqsplit{TestF05TEST13SubmitDraftReportsExpenseDateField} & pass (1 run) \\
F05-TEST-14 & F05-AC-10 & Integration (module) & \path{api/internal/claim/drafts_test.go}\newline \hspace*{1em}\seqsplit{TestF05TEST14DraftSubmissionUsesOwnersCurrentReviewer} & pass (1 run) \\
F05-TEST-15 & F05-AC-11 & Integration (module) & \path{api/internal/claim/drafts_test.go}\newline \hspace*{1em}\seqsplit{TestF05TEST15DraftSubmissionRequiresOwnersCurrentReviewer} & pass (1 run) \\
F05-TEST-16 & F05-AC-12 & Integration & \path{api/claims_drafts_test.go}\newline \hspace*{1em}\seqsplit{TestF05TEST16OnlyDraftOrReturnedCanBeSubmitted} & pass (1 run) \\
F05-TEST-17 & F05-AC-12 & Integration (module) & \path{api/internal/claim/drafts_test.go}\newline \hspace*{1em}\seqsplit{TestF05TEST17ConcurrentDraftSubmissionsOneWins} & pass (1 run) \\
F05-TEST-18 & F05-AC-13 & Integration (module) & \path{api/internal/claim/drafts_test.go}\newline \hspace*{1em}\seqsplit{TestF05TEST18HistoryFailureRollsBackDraftSubmission} & pass (1 run) \\
F05-TEST-19 & F05-AC-14 & Integration & \path{api/claims_drafts_test.go}\newline \hspace*{1em}\seqsplit{TestF05TEST19DeleteRequiresOwnerAndDraftStatus} & pass (1 run) \\
F05-TEST-20 & F05-AC-15 & Integration & \path{api/claims_drafts_test.go}\newline \hspace*{1em}\seqsplit{TestF05TEST20HistoryPreventsDirectClaimDeletion} & pass (1 run) \\
F05-TEST-21 & F05-AC-16 & Component & \path{web/src/features/claims/__test__/draft-workflow.test.tsx}\newline \hspace*{1em}F05-TEST-21: keeps input on a draft save error, disables both actions while saving, and opens the draft & pass (Vitest, CI at 3a91ea4) \\
F05-TEST-22 & F05-AC-17 & Component & \path{web/src/features/claims/__test__/draft-workflow.test.tsx}\newline \hspace*{1em}F05-TEST-22: shows draft placeholders and hides actions absent from policy\newline \hspace*{1em}F05-TEST-22: shows each invalid line, then submits with the version and refetches\newline \hspace*{1em}F05-TEST-22: reloads on conflict and confirms before deleting\newline \hspace*{1em}F05-TEST-22: shows a claim-wide validation message without an unavailable edit action\newline \hspace*{1em}F05-TEST-22: reloads after a rejected transition and shows \$message\newline \hspace*{1em}F05-TEST-22: reports a failed submission when \$name\newline \hspace*{1em}F05-TEST-22: closes the delete confirmation when \$name & pass (Vitest, CI at 3a91ea4) \\
F05-TEST-23 & F05-AC-18 & Component & \path{web/src/features/claims/__test__/draft-workflow.test.tsx}\newline \hspace*{1em}F05-TEST-23: prefills the form, preserves invalid input, and reloads on conflict\newline \hspace*{1em}F05-TEST-23: declines the edit route without the edit action\newline \hspace*{1em}F05-TEST-23: shows the access result when saving returns \$status\newline \hspace*{1em}F05-TEST-23: returns to detail after a rejected edit and shows \$message\newline \hspace*{1em}F05-TEST-23: shows no form error while a conflicting draft reloads\newline \hspace*{1em}F05-TEST-23: retries a failed conflict reload and opens the latest draft & pass (Vitest, CI at 3a91ea4) \\
F05-TEST-24 & F05-AC-19 & Component & \path{web/src/features/claims/__test__/draft-workflow.test.tsx}\newline \hspace*{1em}F05-TEST-24: lists a draft with a neutral badge and a dash in Submitted & pass (Vitest, CI at 3a91ea4) \\
F05-TEST-25 & F05-AC-20 & E2E & \path{web/e2e/f05-drafts.spec.ts}\newline \hspace*{1em}F05-TEST-25: save, edit, submit, review and delete drafts & pass (E2E on a local stack at 3a91ea4, not in CI) \\
F05-TEST-26 & F04-AC-19 (seed) & Integration & \path{api/seed_test.go}\newline \hspace*{1em}\seqsplit{TestF04TEST22SeedAddsSampleClaimsOnce} & pass (1 run) \\
F06-TEST-01 & F06-AC-01 & Integration & \path{api/claims_return_test.go}\newline \hspace*{1em}\seqsplit{TestF06TEST01ReviewerReturnsASubmittedClaim} & pass (1 run) \\
F06-TEST-02 & F06-AC-02 & Integration (module) & \path{api/internal/claim/return_test.go}\newline \hspace*{1em}\seqsplit{TestF06TEST02ReturnReasonFollowsTheRejectRuleBR04} & pass (1 run) \\
F06-TEST-03 & F06-AC-02 & Integration & \path{api/claims_return_test.go}\newline \hspace*{1em}\seqsplit{TestF06TEST03ReturnWithoutAReasonIs422} & pass (1 run) \\
F06-TEST-04 & F06-AC-03 & Integration & \path{api/claims_return_test.go}\newline \hspace*{1em}\seqsplit{TestF06TEST04OnlyTheAssignedReviewerCanReturn} & pass (1 run) \\
F06-TEST-05 & F06-AC-04 & Integration & \path{api/claims_return_test.go}\newline \hspace*{1em}\seqsplit{TestF06TEST05OnlyASubmittedClaimCanBeReturned} & pass (1 run) \\
F06-TEST-06 & F06-AC-05 & Unit & \path{api/internal/claim/policy_test.go}\newline \hspace*{1em}\seqsplit{TestF06TEST06ReturnIsAllowedBelowTheReturnLimit} & pass (1 run) \\
F06-TEST-07 & F06-AC-05 & Integration (module) & \path{api/claims_return_test.go}\newline \hspace*{1em}\seqsplit{TestF06TEST07ReturnLimitIs409AndApproveStillWorks}\newline \path{api/internal/claim/return_test.go}\newline \hspace*{1em}\seqsplit{TestF06TEST07FourthReturnIsRefusedAndApproveStillWorks} & pass (2 run) \\
F06-TEST-08 & F06-AC-06 & Integration & \path{api/claims_return_test.go}\newline \hspace*{1em}\seqsplit{TestF06TEST08ReturnedClaimLeavesTheQueueUntilResubmitted} & pass (1 run) \\
F06-TEST-09 & F06-AC-07 & Integration & \path{api/claims_resubmit_test.go}\newline \hspace*{1em}\seqsplit{TestF06TEST09OwnerResubmitsToTheSameReviewer} & pass (1 run) \\
F06-TEST-10 & F06-AC-08 & Integration (module) & \path{api/internal/claim/resubmit_test.go}\newline \hspace*{1em}\seqsplit{TestF06TEST10ResubmissionCountsNinetyDaysFromTheFirstSubmission} & pass (1 run) \\
F06-TEST-11 & F06-AC-09 & Integration & \path{api/claims_resubmit_test.go}\newline \hspace*{1em}\seqsplit{TestF06TEST11OnlyTheOwnerEditsAndResubmitsAReturnedClaim} & pass (1 run) \\
F06-TEST-12 & F06-AC-10 & Integration (module) & \path{api/internal/claim/resubmit_test.go}\newline \hspace*{1em}\seqsplit{TestF06TEST12HistoryFailureRollsBackTheResubmission}\newline \path{api/internal/claim/return_test.go}\newline \hspace*{1em}\seqsplit{TestF06TEST12HistoryFailureRollsBackTheReturn} & pass (2 run) \\
F06-TEST-13 & F06-AC-11 & Component & \path{web/src/features/reviews/__test__/return-panel.test.tsx}\newline \hspace*{1em}F06-TEST-13: shows Return with the number of returns used\newline \hspace*{1em}F06-TEST-13: shows no Return button when return is not allowed\newline \hspace*{1em}F06-TEST-13: with no returns left it says to approve or reject\newline \hspace*{1em}F06-TEST-13: the dialog counts the reason, and a 422 puts the message under the field with the text kept\newline \hspace*{1em}F06-TEST-13: return requires 10 characters, as reject does\newline \hspace*{1em}F06-TEST-13: the buttons are disabled with a spinner while the return is sent\newline \hspace*{1em}F06-TEST-13: a return shows the toast and goes back to the review queue\newline \hspace*{1em}F06-TEST-13: a 409 shows the changed toast, closes the dialog, and reads the claim again & pass (Vitest, CI at 3a91ea4) \\
F06-TEST-14 & F06-AC-12 & Component & \path{web/src/features/claims/__test__/resubmit.test.tsx}\newline \hspace*{1em}F06-TEST-14 (Edit, Resubmit, toast): the owner of a returned claim sees Edit and Resubmit, not Submit claim, and resubmitting shows the toast and refetches\newline \hspace*{1em}F06-TEST-14 (422): lists each broken line as Line n and keeps the claim returned\newline \path{web/src/features/claims/__test__/return-banner.test.tsx}\newline \hspace*{1em}F06-TEST-14 (banner and count; Edit and Resubmit are \#44): shows the latest return reason with who returned it and when, and the number of returns\newline \hspace*{1em}F06-TEST-14: shows the banner without a reason when the return entry has none & pass (Vitest, CI at 3a91ea4) \\
F06-TEST-15 & F06-AC-13 & E2E & \path{web/e2e/f06-return-resubmit.spec.ts}\newline \hspace*{1em}F06-TEST-15: req.anan submits, rev.malee returns with a reason, req.anan edits and resubmits, rev.malee approves, and the history reads Submit, Return, Resubmit, Approve & pass (E2E on a local stack at 3a91ea4, not in CI) \\
F06-TEST-16 & F04-AC-19 (seed) & Integration & \path{api/seed_test.go}\newline \hspace*{1em}\seqsplit{TestF04TEST22SeedAddsSampleClaimsOnce} & pass (1 run) \\
F06-TEST-17 & F06-AC-07 & Integration (module) & \path{api/internal/claim/resubmit_test.go}\newline \hspace*{1em}\seqsplit{TestF06TEST17ResubmissionWorksWithoutAnOwnerReviewer} & pass (1 run) \\
F07-TEST-01 & F07-AC-01 & Integration & \path{api/claims_withdraw_test.go}\newline \hspace*{1em}\seqsplit{TestF07TEST01OwnerWithdrawsASubmittedOrReturnedClaim} & pass (1 run) \\
F07-TEST-02 & F07-AC-02 & Integration & \path{api/claims_withdraw_test.go}\newline \hspace*{1em}\seqsplit{TestF07TEST02WithdrawCommentIsOptionalWithAtMost500Characters} & pass (1 run) \\
F07-TEST-03 & F07-AC-03 & Integration & \path{api/claims_withdraw_test.go}\newline \hspace*{1em}\seqsplit{TestF07TEST03OnlyASubmittedOrReturnedClaimCanBeWithdrawnBR16} & pass (1 run) \\
F07-TEST-04 & F07-AC-04 & Integration & \path{api/claims_withdraw_test.go}\newline \hspace*{1em}\seqsplit{TestF07TEST04OnlyTheOwnerCanWithdraw} & pass (1 run) \\
F07-TEST-05 & F07-AC-04, F07-AC-06 & Unit & \path{api/internal/claim/policy_test.go}\newline \hspace*{1em}\seqsplit{TestF07TEST05WithdrawIsTheOwnersOnSubmittedOrReturned} & pass (1 run) \\
F07-TEST-06 & F07-AC-05 & Integration & \path{api/claims_withdraw_test.go}\newline \hspace*{1em}\seqsplit{TestF07TEST06WithdrawNeedsTheCurrentIfMatch} & pass (1 run) \\
F07-TEST-07 & F07-AC-06 & Integration & \path{api/claims_withdraw_test.go}\newline \hspace*{1em}\seqsplit{TestF07TEST07AWithdrawnClaimCannotBeActedOnAgain} & pass (1 run) \\
F07-TEST-08 & F07-AC-06 & Integration (module) & \path{api/internal/claim/withdraw_test.go}\newline \hspace*{1em}\seqsplit{TestF07TEST08ConcurrentWithdrawAndApproveHaveOneWinner} & pass (1 run) \\
F07-TEST-09 & F07-AC-07 & Integration (module) & \path{api/internal/claim/withdraw_test.go}\newline \hspace*{1em}\seqsplit{TestF07TEST09HistoryFailureRollsBackTheWithdrawal} & pass (1 run) \\
F07-TEST-10 & F07-AC-08 & Integration & \path{api/claims_withdraw_test.go}\newline \hspace*{1em}\seqsplit{TestF07TEST10AWithdrawnClaimLeavesTheReviewQueueButStaysReadable} & pass (1 run) \\
F07-TEST-11 & F07-AC-09 & Component & \path{web/src/features/claims/__test__/withdraw-claim.test.tsx}\newline \hspace*{1em}F07-TEST-11: shows Withdraw only when the policy gives it\newline \hspace*{1em}F07-TEST-11: opens a dialog with the optional comment, a 0/500 counter and a limit of 500 characters\newline \hspace*{1em}F07-TEST-11: disables the confirm button while sending, then shows the toast, the Withdrawn badge and no actions\newline \hspace*{1em}F07-TEST-11: after the withdrawal it shows the toast, the Withdrawn badge and an empty action bar\newline \hspace*{1em}F07-TEST-11: a 409 CONFLICT shows the reload toast and reads the claim again\newline \hspace*{1em}F07-TEST-11: a 409 INVALID\_TRANSITION shows its message and reads the claim again\newline \hspace*{1em}F07-TEST-11: Cancel and Escape close the dialog without sending anything\newline \hspace*{1em}F07-TEST-11: Escape does not close the dialog while the withdrawal is sending\newline \hspace*{1em}F07-TEST-11: a failure that is not a problem document shows the generic toast and closes the dialog & pass (Vitest, CI at 3a91ea4) \\
F07-TEST-12 & F07-AC-10 & E2E & \path{web/e2e/f07-withdraw.spec.ts}\newline \hspace*{1em}F07-TEST-12: req.anan submits and withdraws a claim with a comment, it shows Withdrawn in My Claims, and rev.malee & pass (E2E on a local stack at 3a91ea4, not in CI) \\
F07-TEST-13 & F04-AC-19 (seed) & Integration & \path{api/seed_test.go}\newline \hspace*{1em}\seqsplit{TestF04TEST22SeedAddsSampleClaimsOnce} & pass (1 run) \\
\bottomrule\end{longtable}}

\subsubsection*{Manual and CI-recorded checks (0)}
{\footnotesize\setlength{\tabcolsep}{3pt}
\begin{longtable}{@{}p{1.9cm}p{2.3cm}p{1.8cm}p{10.2cm}@{}}
\toprule \textbf{Planned test} & \textbf{AC} & \textbf{Level} & \textbf{Status} \\ \midrule \endhead
\bottomrule\end{longtable}}

\subsubsection*{Planned but not implemented yet (0)}
{\footnotesize\setlength{\tabcolsep}{3pt}
\begin{longtable}{@{}p{1.9cm}p{4.2cm}p{2.2cm}p{7.9cm}@{}}
\toprule \textbf{Planned test} & \textbf{AC} & \textbf{Level} & \textbf{Result} \\ \midrule \endhead
\bottomrule\end{longtable}}


\noindent The Result column shows the result of each automated test on \texttt{3a91ea4}: Go tests from the \texttt{go test -json} run above, web component tests from CI run 45614889848, E2E tests from the local Playwright run. The seed rows (F05-TEST-26, F06-TEST-16, F07-TEST-13) are found through \path{docs/traceability.md} and are subtests of \texttt{TestF04TEST22SeedAddsSampleClaimsOnce}; the pass shown is that of the top-level test. There are no manual checks and no planned test without an actual test in Sprint 2.

\subsubsection*{Test excerpts}

\textbf{Domain rule: BR-14, a claim is returned at most three times} (F06-TEST-07, F06-AC-05).
\path{apps/api/internal/claim/return_test.go}, commit \texttt{3a91ea4}, lines 98--132.
The setup is the seeded requester and reviewer and the real \texttt{claim.Service} on PostgreSQL; the helper \texttt{returnAndResubmit} returns and resubmits the claim \texttt{MaxReturns} times; the assertions are that the fourth return fails with \texttt{ErrInvalidTransition} for the return limit, that the stored status, return count and version are unchanged, and that Approve still works.

\VerbatimInput[tabsize=4, fontsize=\scriptsize, frame=single, framesep=5pt, rulecolor=\color{black!25}, breaklines=true]{content/data/excerpt-br14.go}

\textbf{Access restriction: only the owner can withdraw} (F07-TEST-04, F07-AC-04, NFR-02).
\path{apps/api/claims_withdraw_test.go}, commit \texttt{3a91ea4}, lines 147--169.
The setup is the real router with the seeded users: \texttt{req.anan} submits a claim; the assigned reviewer \texttt{rev.malee} and another requester \texttt{req.ploy} each send \texttt{POST /claims/\{id\}/withdraw} with the current version. The assertions are the 403 \texttt{FORBIDDEN} problem, that the body leaks none of \texttt{data}, \texttt{claimNo} or the title, and that the stored claim is unchanged.

\VerbatimInput[tabsize=4, fontsize=\scriptsize, frame=single, framesep=5pt, rulecolor=\color{black!25}, breaklines=true]{content/data/excerpt-f07-access.go}

\textbf{Execution results of the two excerpts} (run on \texttt{3a91ea4}):

\VerbatimInput[tabsize=4, fontsize=\scriptsize, frame=single, framesep=5pt, rulecolor=\color{black!25}, breaklines=true]{content/data/excerpt-results.txt}

\subsubsection*{TDD: the failing test first, then the code}

Two pairs, each a test commit that fails and the commit that makes it pass. All four are commits of PR \#108 (50 commits on \path{feat/f06-f07-return-resubmit-withdraw}); the PR was squash-merged into \texttt{sprint-2} as \texttt{aaf17af}, so they stay reachable under the PR's Commits tab and as \texttt{refs/pull/108/head}. Author of all four: \texttt{Kiatisakk}. Output: the same command run on a clean checkout of each commit (Docker Desktop running); container start-up lines removed.

\begin{center}\small
\begin{tabular}{@{}lllp{7.7cm}@{}}
\toprule \textbf{Rule, test} & \textbf{Step, commit} & \textbf{Author date (UTC+7)} & \textbf{Subject} \\ \midrule
BR-14, F06-TEST-07 & Red \texttt{a547ac0} & 2026-11-25 01:13:33 & \path{test: refuse a fourth return of a claim (F06-TEST-07)} \\
 & Green \texttt{71a8e43} & 2026-11-25 01:17:19 & \path{feat: limit a claim to three returns} \\
BR-16, F07-TEST-03 & Red \texttt{53e8827} & 2026-11-25 02:24:32 & \path{test: only a submitted or returned claim can be withdrawn (F07-TEST-03)} \\
 & Green \texttt{38e0165} & 2026-11-25 02:35:27 & \path{feat: add the withdraw operation for a submitted or returned claim (F07-TEST-03)} \\
\bottomrule
\end{tabular}
\end{center}

\textbf{BR-14: why the red test failed.} At \texttt{a547ac0} the commit adds the test and only the names it needs to compile (the constant \texttt{MaxReturns} and the error \texttt{ErrReturnLimit}); the service does not apply the limit. The fourth return therefore reached PostgreSQL, whose \texttt{CHECK (return\_count <= 3)} refused it, and the service returned a database error where the test expects \texttt{ErrInvalidTransition}. At \texttt{71a8e43} the service checks the count before the update and the test passes.

\textbf{Red, \texttt{a547ac0}:}

\VerbatimInput[tabsize=4, fontsize=\scriptsize, frame=single, framesep=5pt, rulecolor=\color{black!25}, breaklines=true]{content/data/tdd-br14-red.txt}

\textbf{Green, \texttt{71a8e43}:}

\VerbatimInput[tabsize=4, fontsize=\scriptsize, frame=single, framesep=5pt, rulecolor=\color{black!25}, breaklines=true]{content/data/tdd-br14-green.txt}

\textbf{BR-16: why the red test failed.} At \texttt{53e8827} only the test exists: the withdraw route is not in the API, so each of the four starting statuses (\texttt{DRAFT}, \texttt{APPROVED}, \texttt{REJECTED}, \texttt{WITHDRAWN}) got 404 \texttt{NOT\_FOUND} where the test expects 409. At \texttt{38e0165} the operation and its lifecycle rule exist, and all four cases return 409 \texttt{INVALID\_TRANSITION} with BR-16.

\textbf{Red, \texttt{53e8827}:}

\VerbatimInput[tabsize=4, fontsize=\scriptsize, frame=single, framesep=5pt, rulecolor=\color{black!25}, breaklines=true]{content/data/tdd-br16-red.txt}

\textbf{Green, \texttt{38e0165}:}

\VerbatimInput[tabsize=4, fontsize=\scriptsize, frame=single, framesep=5pt, rulecolor=\color{black!25}, breaklines=true]{content/data/tdd-br16-green.txt}

\textbf{Other test-first pairs and honest exceptions.} The Test column of the F-05, F-06 and F-07 result tables in \path{docs/traceability.md} (Part D) lists a red and a green commit for every other test that went red first, for example BR-10 (F05-TEST-09, 19: red \texttt{ec2179d}, green \texttt{dd31979}) and BR-07 (F05-TEST-12: red \texttt{4dcf8d4}, green \texttt{06c5fa0}) in PR \#104. Where a test could not go red first the same table says so and the contract accepts it: F06-TEST-03, 04, 06, 08, 09, 10, 11, 12, 17 and F07-TEST-04, 06 to 10 passed when written because earlier green commits already covered them; F06-TEST-15 and F07-TEST-12 were written after the code; the BR-07 resubmission window (F06-TEST-10) has a mutation check instead of a red commit, accepted in F-06 contract revision 0.3 after review on PR \#108.

\partsection{B.5}{Reviewed integration: Issue and PR chains}\label{B:5}

Two completed chains, each from tickets through a feature branch, a teammate-reviewed PR into the Sprint 2 staging branch \texttt{sprint-2}, and release PR \#109 into \texttt{main}. Times are UTC as GitHub records them. Both PRs were squash-merged into \texttt{sprint-2} (one parent). The closing keywords of the PR bodies do not fire against \texttt{sprint-2}, so the tickets were closed by hand after each merge, as both reviews noted.

\subsubsection*{Chain 1: return, resubmit and withdraw (author 1005, reviewer 1045)}

\begin{center}\small
\begin{tabular}{@{}>{\bfseries}p{3.1cm}p{13.2cm}@{}}
\toprule
Issues & \#43 ``F-06/01: Return a submitted claim'', \#44 ``F-06/02: Edit and resubmit a returned claim'', \#45 ``F-07/01: Withdraw a submitted or returned claim''; all assigned to Kiatisakk, opened 2026-11-13, closed 2026-11-25 08:59 to 09:00 \\
Contract IDs & F-06 (F06-AC-01 to 13, F06-TEST-01 to 17) and F-07 (F07-AC-01 to 10, F07-TEST-01 to 13), merged before implementation in PR \#100 (\texttt{dfc5f6f}, 2026-11-24 15:32) \\
Branch & \path{feat/f06-f07-return-resubmit-withdraw}, from \texttt{sprint-2}; 50 commits at merge, red and green pairs per rule (B.4) \\
Pull request & \#108 into \texttt{sprint-2}, opened 2026-11-24 21:16 by Kiatisakk; 44 files, +3520 \textminus202 \\
Reviewer, verdicts & beambeambeam: changes requested 2026-11-25 07:28 (two blocking, three should-fix, minor and nit items, 11 inline comments); author replied to every thread 08:35; approved 08:47 \\
Integration into \texttt{sprint-2} & squash commit \texttt{aaf17af}, merged 2026-11-25 08:49 by beambeambeam \\
Integration into \texttt{main} & release PR \#109 (\texttt{sprint-2} into \texttt{main}), merge commit \texttt{3a91ea4}, merged 2026-11-25 12:33 by beambeambeam; tag \texttt{v0.1.1.0}; Deploy demo run 45615736377 \\
\bottomrule
\end{tabular}
\end{center}

\textbf{Substantive review comment} (beambeambeam, on \path{apps/web/src/components/status-badge.tsx}): ``\textbf{Blocking:} \texttt{RETURNED} falls back to \texttt{neutral} here, because \texttt{variants} only colours SUBMITTED, APPROVED and REJECTED. F-06 \S2.3 says\ldots'' (the Returned badge is amber). The second blocking item asked for screenshots of every new screen in the PR description. Other findings: the queue-invalidation test (F06-TEST-13) did not prove the invalidation; the withdraw comment limit \texttt{500} was copied, not read from the contract; the return limit \texttt{3} lived in three places; BR-07 for F06-TEST-10 had no red commit.
\textbf{Author's answers} (Kiatisakk, 08:35): the badge is \texttt{yellow} (the design system has no amber), red \texttt{5a844a5}, green \texttt{3d7091e}; the reviewer's mutation check was reproduced, the old test stayed green without the invalidation, and \texttt{f275752} fixed it; the limit is read from the contract in \texttt{5540b51}; the Go message is built from \texttt{MaxReturns} in \texttt{b06c4ec}; the BR-07 evidence rule was proposed as F-06 contract revision 0.3 (\texttt{68d607c}).
\textbf{Approval} (08:47), and the reviewer's closing reply (08:49) accepting the 0.3 note: a mutation check is enough evidence for behaviour that already existed.

\evidencehalf{figures/github/pr108-header.png}{PR \#108, \texttt{feat/f06-f07-return-resubmit-withdraw} into \texttt{sprint-2}: title, state, branch pair.}
\evidencehalf{figures/github/pr108-changes-requested.png}{PR \#108: beambeambeam's review, changes requested, ranked findings.}
\evidencehalf{figures/github/pr108-thread-returned-badge.png}{PR \#108: the blocking finding that the Returned badge is grey, and the author's reply with the red and green commits.}
\evidencehalf{figures/github/pr108-approved.png}{PR \#108: approval after the fixes.}

\subsubsection*{Chain 2: drafts (author 1045, reviewer 1005)}

\begin{center}\small
\begin{tabular}{@{}>{\bfseries}p{3.1cm}p{13.2cm}@{}}
\toprule
Issues & \#39 ``F-05/01: Save and view drafts'', \#40 ``F-05/02: Edit a draft'', \#41 ``F-05/03: Submit a draft'', \#42 ``F-05/04: Delete a draft''; opened 2026-11-13, closed 2026-11-24 17:08 \\
Contract IDs & F-05 (F05-AC-01 to 20, F05-TEST-01 to 26), merged in PR \#100 (\texttt{dfc5f6f}, 2026-11-24 15:32) \\
Branch & \path{feat/implement-issue-39}, from \texttt{sprint-2}; 14 commits at merge, red and green pairs for save, edit and delete, and the expense window at submission \\
Pull request & \#104 into \texttt{sprint-2}, opened 2026-11-24 15:59 by beambeambeam; 40 files, +4139 \textminus502; body ``Closes \#39'' to ``\#42'' \\
Reviewer, verdicts & Kiatisakk: changes requested 16:53 (one blocking item, three non-blocking); author replied 17:01; re-reviewed and approved 17:07 \\
Integration into \texttt{sprint-2} & squash commit \texttt{35052ae}, merged 2026-11-24 17:07 by Kiatisakk \\
Integration into \texttt{main} & release PR \#109, merge commit \texttt{3a91ea4}, 2026-11-25 12:33; tag \texttt{v0.1.1.0} \\
\bottomrule
\end{tabular}
\end{center}

\textbf{Substantive review comment} (Kiatisakk, on \path{docs/traceability.md}): the recorded evidence was older than the head. The Go race suite was recorded on \texttt{e26fb74} and the E2E on \texttt{d5e047e}, but the head \texttt{29c9011} changes production web code after both (\path{claim-load-failure.tsx}, \path{claim-owner-actions.tsx}, \path{edit-claim-page.tsx}). Non-blocking items: a \texttt{test:} commit that changed four non-test files; closing the tickets by hand; a manual merge with the traceability edits of \#101 and \#105.
\textbf{Author's answer} (beambeambeam, 17:01): reran the Go race suite, the web coverage suite and all six Playwright specs on \texttt{29c9011} and recorded them in \texttt{b3cecd1}.
\textbf{Approval} (17:07): ``\texttt{b3cecd1} differs from \texttt{29c9011} only in \path{docs/traceability.md}, so the tested code is the head's code.''

\evidencehalf{figures/github/pr104-header.png}{PR \#104, \texttt{feat/implement-issue-39} into \texttt{sprint-2}.}
\evidencehalf{figures/github/pr104-changes-requested.png}{PR \#104: Kiatisakk's review, changes requested: the recorded evidence is older than the head.}
\evidencehalf{figures/github/pr104-thread-stale-evidence.png}{PR \#104: the finding on the traceability note, and the author's reply with the fix.}
\evidencehalf{figures/github/pr104-approved.png}{PR \#104: approval after re-review, with the check that the tested code is the head's code.}

\subsubsection*{History sequence: contracts before implementation, then both chains to \texttt{main}}

Commands: \texttt{git log -1 --format="\%h  \%ad  \%an  \%s" --date=format:"\%Y-\%m-\%d \%H:\%M \%z" <commit>}, one per commit; \texttt{bbb} is the GitHub account name of beambeambeam.

\begin{shell}
dfc5f6f  2026-11-24 22:32 +0700  bbb        docs: add the Sprint 2 engineering contracts (#100)
35052ae  2026-11-25 00:07 +0700  bbb        feat: add draft claim workflow (#104)
aaf17af  2026-11-25 15:49 +0700  Kiatisakk  feat: return, resubmit and withdraw a claim (#108)
3a91ea4  2026-11-25 19:33 +0700  bbb        chore: release Sprint 2 (#109)   (merge commit on main)
\end{shell}

\begin{center}\small
\begin{tabular}{@{}p{7.2cm}lll@{}}
\toprule \textbf{Step} & \textbf{PR} & \textbf{Commit} & \textbf{Merged (UTC)} \\ \midrule
Sprint 2 contracts and test design (F-05, F-06, F-07) & \#100 & \texttt{dfc5f6f} & 2026-11-24 15:32 \\
Drafts, F-05/01 to 04 (\#39 to \#42) & \#104 & \texttt{35052ae} & 2026-11-24 17:07 \\
Return, resubmit, withdraw, F-06/01, 02 and F-07/01 (\#43 to \#45) & \#108 & \texttt{aaf17af} & 2026-11-25 08:49 \\
Release \texttt{sprint-2} into \texttt{main} & \#109 & \texttt{3a91ea4} (merge) & 2026-11-25 12:33 \\
\bottomrule
\end{tabular}
\end{center}

The contract PR \#100 merged at 22:32 (UTC+7) on 24 November; the first red commit of PR \#104 (\texttt{fad37c0}) is dated 22:36 and the first red commit of PR \#108 (\texttt{93c0d9c}, F06-TEST-01) 00:41 on 25 November. Both implementation chains therefore start after the contracts were merged.

\textbf{Release PR \#109.} Opened 2026-11-25 10:10 by Kiatisakk (17 commits, 67 files); beambeambeam commented on 2026-11-25 10:17 (the merge resolution of \texttt{main} into \texttt{sprint-2}, and stale counts and run IDs in the evidence), the author fixed each in \texttt{bd7c196}, \texttt{c96c609}, \texttt{cb79064} and \texttt{3a9de41}, and he approved the head \texttt{cb79064} at 12:25 after checking the merge resolution and that CI passed on that head (\texttt{code\_quality}, \texttt{test}, \texttt{api}, \texttt{build}, \texttt{image}). Merged with a merge commit at 12:33 so \texttt{main} keeps each feature's subject. Tag \texttt{v0.1.1.0} followed, and the Deploy demo workflow deployed it (Part E).

\evidencehalf{figures/github/pr109-header.png}{Release PR \#109: \texttt{sprint-2} into \texttt{main}, merged with a merge commit.}

Two completed chains, each from tickets through a feature branch, a teammate-reviewed PR into the Sprint 2 staging branch \texttt{sprint-2}, and release PR \#109 into \texttt{main}. Times are UTC as GitHub records them. Both PRs were squash-merged into \texttt{sprint-2} (one parent). The closing keywords of the PR bodies do not fire against \texttt{sprint-2}, so the tickets were closed by hand after each merge, as both reviews noted.

\subsubsection*{Chain 1: return, resubmit and withdraw (author 1005, reviewer 1045)}

\begin{center}\small
\begin{tabular}{@{}>{\bfseries}p{3.1cm}p{13.2cm}@{}}
\toprule
Issues & \#43 ``F-06/01: Return a submitted claim'', \#44 ``F-06/02: Edit and resubmit a returned claim'', \#45 ``F-07/01: Withdraw a submitted or returned claim''; all assigned to Kiatisakk, opened 2026-11-13, closed 2026-11-25 08:59 to 09:00 \\
Contract IDs & F-06 (F06-AC-01 to 13, F06-TEST-01 to 17) and F-07 (F07-AC-01 to 10, F07-TEST-01 to 13), merged before implementation in PR \#100 (\texttt{dfc5f6f}, 2026-11-24 15:32) \\
Branch & \path{feat/f06-f07-return-resubmit-withdraw}, from \texttt{sprint-2}; 50 commits at merge, red and green pairs per rule (B.4) \\
Pull request & \#108 into \texttt{sprint-2}, opened 2026-11-24 21:16 by Kiatisakk; 44 files, +3520 \textminus202 \\
Reviewer, verdicts & beambeambeam: changes requested 2026-11-25 07:28 (two blocking, three should-fix, minor and nit items, 11 inline comments); author replied to every thread 08:35; approved 08:47 \\
Integration into \texttt{sprint-2} & squash commit \texttt{aaf17af}, merged 2026-11-25 08:49 by beambeambeam \\
Integration into \texttt{main} & release PR \#109 (\texttt{sprint-2} into \texttt{main}), merge commit \texttt{3a91ea4}, merged 2026-11-25 12:33 by beambeambeam; tag \texttt{v0.1.1.0}; Deploy demo run 45615736377 \\
\bottomrule
\end{tabular}
\end{center}

\textbf{Substantive review comment} (beambeambeam, on \path{apps/web/src/components/status-badge.tsx}): ``\textbf{Blocking:} \texttt{RETURNED} falls back to \texttt{neutral} here, because \texttt{variants} only colours SUBMITTED, APPROVED and REJECTED. F-06 \S2.3 says\ldots'' (the Returned badge is amber). The second blocking item asked for screenshots of every new screen in the PR description. Other findings: the queue-invalidation test (F06-TEST-13) did not prove the invalidation; the withdraw comment limit \texttt{500} was copied, not read from the contract; the return limit \texttt{3} lived in three places; BR-07 for F06-TEST-10 had no red commit.
\textbf{Author's answers} (Kiatisakk, 08:35): the badge is \texttt{yellow} (the design system has no amber), red \texttt{5a844a5}, green \texttt{3d7091e}; the reviewer's mutation check was reproduced, the old test stayed green without the invalidation, and \texttt{f275752} fixed it; the limit is read from the contract in \texttt{5540b51}; the Go message is built from \texttt{MaxReturns} in \texttt{b06c4ec}; the BR-07 evidence rule was proposed as F-06 contract revision 0.3 (\texttt{68d607c}).
\textbf{Approval} (08:47), and the reviewer's closing reply (08:49) accepting the 0.3 note: a mutation check is enough evidence for behaviour that already existed.

\evidencehalf{figures/github/pr108-header.png}{PR \#108, \texttt{feat/f06-f07-return-resubmit-withdraw} into \texttt{sprint-2}: title, state, branch pair.}
\evidencehalf{figures/github/pr108-changes-requested.png}{PR \#108: beambeambeam's review, changes requested, ranked findings.}
\evidencehalf{figures/github/pr108-thread-returned-badge.png}{PR \#108: the blocking finding that the Returned badge is grey, and the author's reply with the red and green commits.}
\evidencehalf{figures/github/pr108-approved.png}{PR \#108: approval after the fixes.}

\subsubsection*{Chain 2: drafts (author 1045, reviewer 1005)}

\begin{center}\small
\begin{tabular}{@{}>{\bfseries}p{3.1cm}p{13.2cm}@{}}
\toprule
Issues & \#39 ``F-05/01: Save and view drafts'', \#40 ``F-05/02: Edit a draft'', \#41 ``F-05/03: Submit a draft'', \#42 ``F-05/04: Delete a draft''; opened 2026-11-13, closed 2026-11-24 17:08 \\
Contract IDs & F-05 (F05-AC-01 to 20, F05-TEST-01 to 26), merged in PR \#100 (\texttt{dfc5f6f}, 2026-11-24 15:32) \\
Branch & \path{feat/implement-issue-39}, from \texttt{sprint-2}; 14 commits at merge, red and green pairs for save, edit and delete, and the expense window at submission \\
Pull request & \#104 into \texttt{sprint-2}, opened 2026-11-24 15:59 by beambeambeam; 40 files, +4139 \textminus502; body ``Closes \#39'' to ``\#42'' \\
Reviewer, verdicts & Kiatisakk: changes requested 16:53 (one blocking item, three non-blocking); author replied 17:01; re-reviewed and approved 17:07 \\
Integration into \texttt{sprint-2} & squash commit \texttt{35052ae}, merged 2026-11-24 17:07 by Kiatisakk \\
Integration into \texttt{main} & release PR \#109, merge commit \texttt{3a91ea4}, 2026-11-25 12:33; tag \texttt{v0.1.1.0} \\
\bottomrule
\end{tabular}
\end{center}

\textbf{Substantive review comment} (Kiatisakk, on \path{docs/traceability.md}): the recorded evidence was older than the head. The Go race suite was recorded on \texttt{e26fb74} and the E2E on \texttt{d5e047e}, but the head \texttt{29c9011} changes production web code after both (\path{claim-load-failure.tsx}, \path{claim-owner-actions.tsx}, \path{edit-claim-page.tsx}). Non-blocking items: a \texttt{test:} commit that changed four non-test files; closing the tickets by hand; a manual merge with the traceability edits of \#101 and \#105.
\textbf{Author's answer} (beambeambeam, 17:01): reran the Go race suite, the web coverage suite and all six Playwright specs on \texttt{29c9011} and recorded them in \texttt{b3cecd1}.
\textbf{Approval} (17:07): ``\texttt{b3cecd1} differs from \texttt{29c9011} only in \path{docs/traceability.md}, so the tested code is the head's code.''

\evidencehalf{figures/github/pr104-header.png}{PR \#104, \texttt{feat/implement-issue-39} into \texttt{sprint-2}.}
\evidencehalf{figures/github/pr104-changes-requested.png}{PR \#104: Kiatisakk's review, changes requested: the recorded evidence is older than the head.}
\evidencehalf{figures/github/pr104-thread-stale-evidence.png}{PR \#104: the finding on the traceability note, and the author's reply with the fix.}
\evidencehalf{figures/github/pr104-approved.png}{PR \#104: approval after re-review, with the check that the tested code is the head's code.}

\subsubsection*{History sequence: contracts before implementation, then both chains to \texttt{main}}

Commands: \texttt{git log -1 --format="\%h  \%ad  \%an  \%s" --date=format:"\%Y-\%m-\%d \%H:\%M \%z" <commit>}, one per commit; \texttt{bbb} is the GitHub account name of beambeambeam.

\begin{shell}
dfc5f6f  2026-11-24 22:32 +0700  bbb        docs: add the Sprint 2 engineering contracts (#100)
35052ae  2026-11-25 00:07 +0700  bbb        feat: add draft claim workflow (#104)
aaf17af  2026-11-25 15:49 +0700  Kiatisakk  feat: return, resubmit and withdraw a claim (#108)
3a91ea4  2026-11-25 19:33 +0700  bbb        chore: release Sprint 2 (#109)   (merge commit on main)
\end{shell}

\begin{center}\small
\begin{tabular}{@{}p{7.2cm}lll@{}}
\toprule \textbf{Step} & \textbf{PR} & \textbf{Commit} & \textbf{Merged (UTC)} \\ \midrule
Sprint 2 contracts and test design (F-05, F-06, F-07) & \#100 & \texttt{dfc5f6f} & 2026-11-24 15:32 \\
Drafts, F-05/01 to 04 (\#39 to \#42) & \#104 & \texttt{35052ae} & 2026-11-24 17:07 \\
Return, resubmit, withdraw, F-06/01, 02 and F-07/01 (\#43 to \#45) & \#108 & \texttt{aaf17af} & 2026-11-25 08:49 \\
Release \texttt{sprint-2} into \texttt{main} & \#109 & \texttt{3a91ea4} (merge) & 2026-11-25 12:33 \\
\bottomrule
\end{tabular}
\end{center}

The contract PR \#100 merged on 2026-11-24 at 15:32 UTC, 3.5 minutes before the first red commit of PR \#104 (\texttt{fad37c0}, 15:36 UTC). The first commit of PR \#108 is a test commit too, \texttt{93c0d9c} (``test: return a submitted claim (F06-TEST-01)'', 2026-11-24 17:41 UTC): 2 hours 9 minutes after the merge of \#100, and 34 minutes after the merge of \#104. In both chains the contract was merged before the first test was written, and the test commit comes before the code commit it drives (\texttt{630fca4}, 18:01 UTC, for \texttt{93c0d9c}).

\textbf{Release PR \#109.} Opened 2026-11-25 10:10 by Kiatisakk (17 commits, 67 files); beambeambeam commented on 2026-11-25 10:17 (the merge resolution of \texttt{main} into \texttt{sprint-2}, and stale counts and run IDs in the evidence), the author fixed each in \texttt{bd7c196}, \texttt{c96c609}, \texttt{cb79064} and \texttt{3a9de41}, and he approved the head \texttt{cb79064} at 12:25 after checking the merge resolution and that CI passed on that head (\texttt{code\_quality}, \texttt{test}, \texttt{api}, \texttt{build}, \texttt{image}). Merged with a merge commit at 12:33 so \texttt{main} keeps each feature's subject. Tag \texttt{v0.1.1.0} followed, and the Deploy demo workflow deployed it (Part E).

\evidencehalf{figures/github/pr109-header.png}{Release PR \#109: \texttt{sprint-2} into \texttt{main}, merged with a merge commit.}

\reportpartend{B}

\reportpart{C}{Sprint 3 contracts and working progress}
\partsection{C.1}{Sprint 3 contracts}\label{C:1}

\todo{beam: write C.1. The Sprint 3 contracts are Beam's work and not merged yet (PR \#111 targets \texttt{sprint-3} and is open). Render F-08, F-09 and F-10 and their DoD here through \path{report.config.ts} once the contracts are on a branch the report can read}

\begin{center}\small
\begin{tabular}{@{}lll@{}}
\toprule \textbf{Feature} & \textbf{Contract} & \textbf{Tickets} \\ \midrule
F-08 Finance Queue and final decisions & \todo{beam: revision} & \#46, \#47 \\
F-09 All-claims view with filters and search & \todo{beam: revision} & \#48 \\
F-10 Decision summary and load dataset (NFR-07) & \todo{beam: revision} & \#49, \#50 \\
\bottomrule
\end{tabular}
\end{center}

\partsection{C.2}{Implemented progress with tests}\label{C:2}

\todo{beam: describe what of F-08, F-09 and F-10 is implemented at the submission commit, with the tests that pass (IDs, level, result) and a short demonstration; write ``nothing implemented'' for any feature without code}

\partsection{C.3}{Feature status}\label{C:3}

\begin{center}\small
\begin{tabular}{@{}lccccc@{}}
\toprule \textbf{Feature} & \textbf{Planned} & \textbf{In progress} & \textbf{Implemented} & \textbf{Tested} & \textbf{Accepted} \\ \midrule
F-08 & \todo{beam: mark} & & & & \\
F-09 & \todo{beam: mark} & & & & \\
F-10 & \todo{beam: mark} & & & & \\
\bottomrule
\end{tabular}
\end{center}

\partsection{C.4}{Remaining work}\label{C:4}

\begin{center}\small
\begin{tabular}{@{}p{3.2cm}p{2.5cm}p{5.8cm}p{4.6cm}@{}}
\toprule \textbf{Feature} & \textbf{Status} & \textbf{Remaining action} & \textbf{Owner and target date} \\ \midrule
F-08 Finance Queue & \todo{beam: status} & \todo{beam: action} & \todo{beam: owner, date} \\
F-09 All-claims view & \todo{beam: status} & \todo{beam: action} & \todo{beam: owner, date} \\
F-10 Decision summary, NFR-07 run & \todo{beam: status} & \todo{beam: action} & \todo{beam: owner, date} \\
\bottomrule
\end{tabular}
\end{center}

\reportpartend{C}

\reportpart{D}{Updated requirements and test traceability}
% Generated by scripts/gen-part-lists.ts. Do not edit.

\partsection{Software Requirements Specification}{}
\docsource{docs/srs.md}{main}{3a91ea4}
\evidencepart{figures/docs/srs-p1.png}{\detokenize{docs/srs.md, piece 1 of 8}}
\evidencepart{figures/docs/srs-p2.png}{\detokenize{docs/srs.md, piece 2 of 8}}
\evidencepart{figures/docs/srs-p3.png}{\detokenize{docs/srs.md, piece 3 of 8}}
\evidencepart{figures/docs/srs-p4.png}{\detokenize{docs/srs.md, piece 4 of 8}}
\evidencepart{figures/docs/srs-p5.png}{\detokenize{docs/srs.md, piece 5 of 8}}
\evidencepart{figures/docs/srs-p6.png}{\detokenize{docs/srs.md, piece 6 of 8}}
\evidencepart{figures/docs/srs-p7.png}{\detokenize{docs/srs.md, piece 7 of 8}}
\evidencepart{figures/docs/srs-p8.png}{\detokenize{docs/srs.md, piece 8 of 8}}

\partsection{Feature inventory}{}
\docsource{docs/features.md}{main}{3a91ea4}
\evidencepart{figures/docs/features-p1.png}{\detokenize{docs/features.md, piece 1 of 4}}
\evidencepart{figures/docs/features-p2.png}{\detokenize{docs/features.md, piece 2 of 4}}
\evidencepart{figures/docs/features-p3.png}{\detokenize{docs/features.md, piece 3 of 4}}
\evidencepart{figures/docs/features-p4.png}{\detokenize{docs/features.md, piece 4 of 4}}
\noindent\textit{The 1 diagram of \path{docs/features.md}, repeated at full size from the same commit (3a91ea4):}
\diagramfull{figures/diagrams/features-1.png}{\detokenize{docs/features.md}, diagram 1 of 1}

% Generated by scripts/gen-part-lists.ts. Do not edit.

\partsection{System-level SDS}{}
\docsource{docs/system-sds.md}{main}{3a91ea4}
\evidencepart{figures/docs/system-sds-p1.png}{\detokenize{docs/system-sds.md, piece 1 of 23}}
\evidencepart{figures/docs/system-sds-p2.png}{\detokenize{docs/system-sds.md, piece 2 of 23}}
\evidencepart{figures/docs/system-sds-p3.png}{\detokenize{docs/system-sds.md, piece 3 of 23}}
\evidencepart{figures/docs/system-sds-p4.png}{\detokenize{docs/system-sds.md, piece 4 of 23}}
\evidencepart{figures/docs/system-sds-p5.png}{\detokenize{docs/system-sds.md, piece 5 of 23}}
\evidencepart{figures/docs/system-sds-p6.png}{\detokenize{docs/system-sds.md, piece 6 of 23}}
\evidencepart{figures/docs/system-sds-p7.png}{\detokenize{docs/system-sds.md, piece 7 of 23}}
\evidencepart{figures/docs/system-sds-p8.png}{\detokenize{docs/system-sds.md, piece 8 of 23}}
\evidencepart{figures/docs/system-sds-p9.png}{\detokenize{docs/system-sds.md, piece 9 of 23}}
\evidencepart{figures/docs/system-sds-p10.png}{\detokenize{docs/system-sds.md, piece 10 of 23}}
\evidencepart{figures/docs/system-sds-p11.png}{\detokenize{docs/system-sds.md, piece 11 of 23}}
\evidencepart{figures/docs/system-sds-p12.png}{\detokenize{docs/system-sds.md, piece 12 of 23}}
\evidencepart{figures/docs/system-sds-p13.png}{\detokenize{docs/system-sds.md, piece 13 of 23}}
\evidencepart{figures/docs/system-sds-p14.png}{\detokenize{docs/system-sds.md, piece 14 of 23}}
\evidencepart{figures/docs/system-sds-p15.png}{\detokenize{docs/system-sds.md, piece 15 of 23}}
\evidencepart{figures/docs/system-sds-p16.png}{\detokenize{docs/system-sds.md, piece 16 of 23}}
\evidencepart{figures/docs/system-sds-p17.png}{\detokenize{docs/system-sds.md, piece 17 of 23}}
\evidencepart{figures/docs/system-sds-p18.png}{\detokenize{docs/system-sds.md, piece 18 of 23}}
\evidencepart{figures/docs/system-sds-p19.png}{\detokenize{docs/system-sds.md, piece 19 of 23}}
\evidencepart{figures/docs/system-sds-p20.png}{\detokenize{docs/system-sds.md, piece 20 of 23}}
\evidencepart{figures/docs/system-sds-p21.png}{\detokenize{docs/system-sds.md, piece 21 of 23}}
\evidencepart{figures/docs/system-sds-p22.png}{\detokenize{docs/system-sds.md, piece 22 of 23}}
\evidencepart{figures/docs/system-sds-p23.png}{\detokenize{docs/system-sds.md, piece 23 of 23}}
\noindent\textit{The 3 diagrams of \path{docs/system-sds.md}, repeated at full size from the same commit (3a91ea4):}
\diagramwide{figures/diagrams/system-sds-1.png}{\detokenize{docs/system-sds.md}, diagram 1 of 3}
\diagramfull{figures/diagrams/system-sds-2.png}{\detokenize{docs/system-sds.md}, diagram 2 of 3}
\diagramwide{figures/diagrams/system-sds-3.png}{\detokenize{docs/system-sds.md}, diagram 3 of 3}

\partsection{ADR 0001}{}
\docsource{docs/adr/0001-finance-officer-reviews-department-head.md}{main}{3a91ea4}
\evidencepart{figures/docs/adr-0001.png}{\detokenize{docs/adr/0001-finance-officer-reviews-department-head.md}}

% Generated by scripts/gen-part-lists.ts. Do not edit.

\partsection{Requirements traceability matrix}{}
\docsource{docs/traceability.md}{main}{3a91ea4}
\evidencepart{figures/docs/traceability-p1.png}{\detokenize{docs/traceability.md, piece 1 of 31}}
\evidencepart{figures/docs/traceability-p2.png}{\detokenize{docs/traceability.md, piece 2 of 31}}
\evidencepart{figures/docs/traceability-p3.png}{\detokenize{docs/traceability.md, piece 3 of 31}}
\evidencepart{figures/docs/traceability-p4.png}{\detokenize{docs/traceability.md, piece 4 of 31}}
\evidencepart{figures/docs/traceability-p5.png}{\detokenize{docs/traceability.md, piece 5 of 31}}
\evidencepart{figures/docs/traceability-p6.png}{\detokenize{docs/traceability.md, piece 6 of 31}}
\evidencepart{figures/docs/traceability-p7.png}{\detokenize{docs/traceability.md, piece 7 of 31}}
\evidencepart{figures/docs/traceability-p8.png}{\detokenize{docs/traceability.md, piece 8 of 31}}
\evidencepart{figures/docs/traceability-p9.png}{\detokenize{docs/traceability.md, piece 9 of 31}}
\evidencepart{figures/docs/traceability-p10.png}{\detokenize{docs/traceability.md, piece 10 of 31}}
\evidencepart{figures/docs/traceability-p11.png}{\detokenize{docs/traceability.md, piece 11 of 31}}
\evidencepart{figures/docs/traceability-p12.png}{\detokenize{docs/traceability.md, piece 12 of 31}}
\evidencepart{figures/docs/traceability-p13.png}{\detokenize{docs/traceability.md, piece 13 of 31}}
\evidencepart{figures/docs/traceability-p14.png}{\detokenize{docs/traceability.md, piece 14 of 31}}
\evidencepart{figures/docs/traceability-p15.png}{\detokenize{docs/traceability.md, piece 15 of 31}}
\evidencepart{figures/docs/traceability-p16.png}{\detokenize{docs/traceability.md, piece 16 of 31}}
\evidencepart{figures/docs/traceability-p17.png}{\detokenize{docs/traceability.md, piece 17 of 31}}
\evidencepart{figures/docs/traceability-p18.png}{\detokenize{docs/traceability.md, piece 18 of 31}}
\evidencepart{figures/docs/traceability-p19.png}{\detokenize{docs/traceability.md, piece 19 of 31}}
\evidencepart{figures/docs/traceability-p20.png}{\detokenize{docs/traceability.md, piece 20 of 31}}
\evidencepart{figures/docs/traceability-p21.png}{\detokenize{docs/traceability.md, piece 21 of 31}}
\evidencepart{figures/docs/traceability-p22.png}{\detokenize{docs/traceability.md, piece 22 of 31}}
\evidencepart{figures/docs/traceability-p23.png}{\detokenize{docs/traceability.md, piece 23 of 31}}
\evidencepart{figures/docs/traceability-p24.png}{\detokenize{docs/traceability.md, piece 24 of 31}}
\evidencepart{figures/docs/traceability-p25.png}{\detokenize{docs/traceability.md, piece 25 of 31}}
\evidencepart{figures/docs/traceability-p26.png}{\detokenize{docs/traceability.md, piece 26 of 31}}
\evidencepart{figures/docs/traceability-p27.png}{\detokenize{docs/traceability.md, piece 27 of 31}}
\evidencepart{figures/docs/traceability-p28.png}{\detokenize{docs/traceability.md, piece 28 of 31}}
\evidencepart{figures/docs/traceability-p29.png}{\detokenize{docs/traceability.md, piece 29 of 31}}
\evidencepart{figures/docs/traceability-p30.png}{\detokenize{docs/traceability.md, piece 30 of 31}}
\evidencepart{figures/docs/traceability-p31.png}{\detokenize{docs/traceability.md, piece 31 of 31}}

\partsection{D.4}{Coverage and gaps}\label{D:4}

From \path{docs/traceability.md} section 4, checked against the SRS (the check date is below). Sprint 2 contracts were checked on 2026-11-23.

\begin{center}\small
\begin{tabular}{@{}lllll@{}}
\toprule \textbf{Type} & \textbf{In SRS} & \textbf{Traced} & \textbf{Mapped to a feature or SR} & \textbf{Uncovered} \\ \midrule
Functional (FR) & 14 & 14 & 14 & 0 \\
Business rules (BR) & 17 & 17 & 17 & 0 \\
Non-functional (NFR) & 10 & 10 & 10 & 0 \\
\textbf{Total} & \textbf{41} & \textbf{41} & \textbf{41} & \textbf{0} \\
\bottomrule
\end{tabular}
\end{center}

Every one of the 43 Sprint 2 AC has at least one of the 56 planned tests, and all 56 have an actual passing test (B.2, B.4). The gaps that remain, stated plainly:

\begin{itemize}\setlength\itemsep{2pt}
\item \textbf{Sprint 3 contracts are on \texttt{sprint-3}, not on \texttt{main}} (F-08, F-09, F-10; PR \#111 merged into \texttt{sprint-3} at \texttt{ef38550}, C.1). The matrix shown here is the one on \texttt{main}: its requirements are mapped to features, but the AC and tests of the Sprint 3 contracts reach \texttt{main} with the final release, and no Sprint 3 test has run.
\item \textbf{NFR-07 is not measured yet.} The timed API run on at least 200 seeded claims is planned for Sprint 3 on the deployed instance.
\item \textbf{E2E specs do not run in CI} (Product DoD 0.3): they run on a local seeded stack before each release (8 of 8 on \texttt{3a91ea4}).
\item \textbf{Web suite local run.} A full local run of the web suite times out under load, so CI is its verdict; one test (F01-TEST-20) fails locally alone and passes in CI.
\item \textbf{Receipt upload} is excluded by the SRS (EX-05); a receipt is a text reference.
\item \textbf{NFR-04 (cloud half), NFR-05 and NFR-10} are verified in Lab 5, not in a sprint test (listed as verified outside a sprint deliverable).
\end{itemize}

\begin{sloppypar}
\textbf{Re-check on 2026-11-26.} \texttt{git diff --stat 3a91ea4 origin/main} restricted to \path{docs/traceability.md}, \path{docs/srs.md} and \path{docs/features.md} is empty (\texttt{origin/main} is \texttt{1ee2563}), and section 4 of \path{docs/traceability.md} on \texttt{main} gives the same figures as the table above: FR 14/14/14/0, BR 17/17/17/0, NFR 10/10/10/0, total 41/41/41/0, ``Uncovered requirements: none''. The file states that its check was made against SRS 0.2 on 2026-11-04; the SRS is 0.3 since L7/01 (A.3), which changed a value of BR-04 and no requirement ID.
\end{sloppypar}

\reportpartend{D}

\reportpart{E}{Integrated release candidate}
\partsection{E.1}{Release candidate: build, test, CI and deploy}\label{E:1}

The candidate is the Sprint 2 release \texttt{v0.1.1.0}, commit \texttt{3a91ea4} on \texttt{main}. It stands in until the Sprint 3 candidate exists. \todo{recapture on candidate: replace tag, commit, CI run, deploy run, revision and image digest with those of the submitted candidate}

\begin{center}\small
\begin{tabular}{@{}>{\bfseries}p{3.6cm}p{12.7cm}@{}}
\toprule
Tag, commit & \texttt{v0.1.1.0}, \texttt{3a91ea4} (merge of release PR \#109), 2026-11-25 19:33 (UTC+7) \\
CI & run 45614889848 on \texttt{3a91ea4}: \texttt{image}, \texttt{code\_quality}, \texttt{test}, \texttt{build}, \texttt{api}, \texttt{release} all \texttt{success} \\
Deploy demo & run 45615736377: Cloud Run revision \texttt{claimflow-00018-kzz}, image digest \texttt{sha256:dd4d3477\ldots}; smoke test: \texttt{/api/v1/health} reports the SHA, \texttt{/api/v1/ready} 200, \texttt{/} serves the web build, an unknown \texttt{/api} route returns a JSON 404 \\
Deployed URL & \url{https://claimflow-5bcxewbqkq-as.a.run.app} \\
Live check & 2026-11-25 about 20:20 (UTC+7): \texttt{GET /api/v1/health} returned version \texttt{3a91ea458862a02fc88384f09ef74b7a502ce762}; \texttt{GET /api/v1/ready} returned 200 \\
\bottomrule
\end{tabular}
\end{center}

\textbf{Tests on the candidate.} Go: 178 top-level tests and 241 subtests, 0 failed, 0 skipped; race run with 100\% coverage (956 of 956). Web: 182 of 182 in CI with 100\% coverage. E2E: 8 of 8 on a local seeded stack (all in B.4).
\textbf{Build.} \texttt{bun run build} and the container image build pass in CI (jobs \texttt{build} and \texttt{image}).

\partsection{E.2}{Setup summary}\label{E:2}

Full steps are in \path{README.md}, \path{apps/api/README.md} and \path{apps/web/README.md}; they were rerun on a clean clone during the Sprint 2 release checks (Windows, A.2). In short:

\begin{enumerate}\setlength\itemsep{2pt}
\item \texttt{bun install} at the repository root.
\item Copy each \texttt{.env.example} to \texttt{.env} and set the values (Varlock loads them against \texttt{.env.schema}); \texttt{bun run db:up} starts PostgreSQL 17 in Docker.
\item In \path{apps/api}: \texttt{migrate up}, then \texttt{seed} (7 users, 9 sample claims; a second run adds none), then \texttt{bun run dev}.
\item In \path{apps/web}: \texttt{bun run dev}.
\end{enumerate}

\textbf{Environment variable names} (values are never in this report or in the repository). API: \texttt{GIN\_MODE}, \texttt{PORT}, \texttt{DATABASE\_URL}, \texttt{COOKIE\_SECURE}, \texttt{BCRYPT\_COST}, \texttt{LOG\_LEVEL}, \texttt{WEB\_DIR}, \texttt{APP\_VERSION}, \texttt{SEED\_USER\_PASSWORD}. Web: the \texttt{VITE\_}-prefixed variables of \path{apps/web/.env.schema}. Report: \texttt{STUDENT\_1\_NAME}, \texttt{STUDENT\_1\_ID}, \texttt{STUDENT\_2\_NAME}, \texttt{STUDENT\_2\_ID}.

\textbf{Deployment.} A merge to \texttt{main} builds the image and runs the Deploy demo workflow to Cloud Run on GCP; credentials are Workload Identity and repository secrets, none in the code. \todo{recapture on candidate: confirm the setup steps on a clean clone of the candidate}

\partsection{E.3}{Readiness review and plan to 13 December}\label{E:3}

\textbf{Readiness of \texttt{v0.1.1.0}.} Ready for Sprint 2 acceptance: all 43 Sprint 2 AC have passing tests, CI and the deploy are green, the 100\% coverage gates hold, and the register (A.2) has no open item. Not ready as a final product: F-08, F-09 and F-10 are not implemented, NFR-07 is not measured, and the E2E specs run locally only. Deliverable 2 (section 2) asks for at least one tested Sprint 3 workflow or rule in the candidate on \texttt{main}; this candidate has none (C.2). The Sprint 3 contracts are merged, and F-08 has started (PR \#113, open).

\textbf{Plan} (the dates of C.4; they slip from \path{docs/sprints/plan.md} section 2.3, which set F-09 and F-10 for 30 November, as C.4 explains):

\begin{center}\small\setlength{\tabcolsep}{4pt}
\begin{tabular}{@{}p{3.0cm}p{8.2cm}p{4.6cm}@{}}
\toprule \textbf{Date} & \textbf{Work} & \textbf{Owner} \\ \midrule
done, 25 Nov & Sprint 3 contracts merged (PR \#111); F-08 started (PR \#113) & \studentTwoName{} (\texttt{@beambeambeam}) \\
by 3 Dec & F-08 merged into \texttt{sprint-3}, tests green & \studentTwoName{} (\texttt{@beambeambeam}) \\
by 6 Dec & F-09 merged & \studentTwoName{} (\texttt{@beambeambeam}) \\
by 8 Dec & F-10 merged, with the load dataset & \studentOneName{} (\texttt{@Kiatisakk}) \\
by 10 Dec & Final regression, UAT re-tests, NFR-07 timed run on the deployed release & both \\
by 13 Dec & Final video, final release to \texttt{main}, deploy, report; due 13 Dec 23:59 on LEB2 & both \\
\bottomrule
\end{tabular}
\end{center}

\reportpartend{E}

\reportback{Back matter: contributions, AI use, and known limitations}
\subsubsection*{Contribution table}

From the merged pull requests and closed issues of \path{beambeambeam/claimflow}, collected by \path{docs/reports/scripts/collect-team.ts}, counting only work from 16 November 2026 (the Sprint 2 period). Dependabot pull requests are not counted.

% Generated by scripts/gen-team-table.ts from team.json. Do not edit.
\begin{center}\small
\begin{tabular}{@{}p{4.8cm}p{5.2cm}p{5.2cm}@{}}
\toprule & \textbf{\studentOneTag{} Kiatisakk} & \textbf{\studentTwoTag{} beambeambeam} \\ \midrule
Merged pull requests authored & 11 & 16 \\
Closed issues assigned & 9 & 0 \\
PRs reviewed (teammate's) & 16 & 10 \\
\quad with an approval & 16 & 10 \\
\quad with changes requested & 10 & 1 \\
Pull requests, numbers & \#64, \#74, \#75, \#76, \#77, \#79, \#80, \#94, \#101, \#108, \#109 & \#72, \#81, \#85, \#88, \#91, \#92, \#93, \#95, \#96, \#97, \#98, \#100, \#104, \#105, \#106, \#107 \\
Issues, numbers & \#4, \#43, \#44, \#45, \#52, \#53, \#54, \#55, \#99 &  \\
\bottomrule\end{tabular}\end{center}

{\footnotesize\begin{longtable}{@{}p{0.9cm}p{7.4cm}p{2.6cm}p{3.9cm}@{}}
\toprule \textbf{PR} & \textbf{Title} & \textbf{Author} & \textbf{Review by the other student} \\ \midrule \endhead
\#64 & docs: add the Sprint 1 report sources & Kiatisakk & none recorded \\
\#72 & docs: add Lab 5 and Lab 6 specification documents & beambeambeam & approved \\
\#74 & feat: one cloud-ready image with the web build and the compose stack & Kiatisakk & approved, commented \\
\#75 & build: add the one-time GCP setup and first deploy script & Kiatisakk & approved, commented \\
\#76 & chore: release Lab 5 image and GCP setup & Kiatisakk & approved, commented, dismissed \\
\#77 & docs: match the ruleset text to the live rulesets & Kiatisakk & approved \\
\#79 & build: deploy main to Cloud Run with a smoke test and rollback & Kiatisakk & approved, commented \\
\#80 & chore: release the Lab 5 deploy pipeline & Kiatisakk & approved \\
\#81 & ci: speed up the API tests and upload the junit reports & beambeambeam & approved \\
\#85 & docs: add the Lab 6 Sprint 1 UAT cases and handoff & beambeambeam & approved, changes requested \\
\#88 & chore: track the agent skills in git and update them & beambeambeam & approved \\
\#91 & fix: prevent a double-click on Submit from creating two claims & beambeambeam & approved, changes requested \\
\#92 & fix: reject claim text that is only whitespace & beambeambeam & approved, changes requested \\
\#93 & chore: release Lab 6 & beambeambeam & approved \\
\#94 & docs: add the Lab 5 report source & Kiatisakk & approved \\
\#95 & docs: add the Lab 6 UAT report & beambeambeam & approved, changes requested \\
\#96 & docs: run E2E specs on a local seeded stack, not in CI & beambeambeam & approved \\
\#97 & chore: Sprint 1 follow-ups & beambeambeam & approved, changes requested \\
\#98 & docs: add the Lab 7 sprint 1 maintenance specification & beambeambeam & approved, changes requested \\
\#100 & docs: add the Sprint 2 engineering contracts & beambeambeam & approved, changes requested \\
\#101 & docs: raise the reject reason minimum to 10 characters (L7/01) & Kiatisakk & approved, commented \\
\#104 & feat: add draft claim workflow & beambeambeam & approved, changes requested \\
\#105 & feat!: raise the reject reason minimum to 10 characters & beambeambeam & approved, changes requested \\
\#106 & chore: release Lab 7 & beambeambeam & approved \\
\#107 & docs: add the Lab 7 maintenance report & beambeambeam & approved, changes requested \\
\#108 & feat: return, resubmit and withdraw a claim & Kiatisakk & approved, changes requested, commented \\
\#109 & chore: release Sprint 2 & Kiatisakk & approved, commented \\
\bottomrule\end{longtable}}


\todo{beam: add one paragraph on your Sprint 2 work (the draft workflow \#104, the contracts \#100, the release review \#109), as in Sprint 1, and confirm the split}

\subsubsection*{AI-use declaration}

\textbf{Student \studentOneId{} (\studentOneName).} Claude Code (Anthropic) was used throughout Sprint 2: implementation of the F-06 and F-07 tickets under test-driven development (the failing test is committed first, B.4), drafting of review comments on the teammate's pull requests, and the capture and generation scripts and drafts of this report. Every result is gated by the repository checks and was read by the student before it was committed, pushed or posted.

\textbf{Student \studentTwoTag{} (\studentTwoName).} \todo{beam: declare your AI use in Sprint 2 (tools, models, tasks, how results were checked)}

\subsubsection*{Known limitations, failed and skipped tests}

\begin{itemize}\setlength\itemsep{2pt}
\item \textbf{E2E specs do not run in CI} (Product DoD 0.3): 8 of 8 pass on a local seeded stack.
\item \textbf{Web tests locally.} A full local web run times out under load; \texttt{F01-TEST-20} (\path{auth-routing.test.tsx}) fails alone on the author's machine and passes in CI. CI is the verdict.
\item \textbf{Receipt files are not supported} (EX-05); issue \#61 was closed without an implementation.
\item \textbf{Sprint 3 is not implemented:} F-08, F-09, F-10 and the NFR-07 timed run (Part C).
\item \textbf{Failed and skipped tests on \texttt{3a91ea4}.} Go: 0 failed, 0 skipped of 178 top-level tests and 241 subtests. Web: 182 of 182 in CI. E2E: 8 of 8 locally. No test was disabled, skipped or weakened to obtain these results.
\end{itemize}

\reportbackend

\end{document}
